\documentclass[a4paper,10pt]{article}
\usepackage[UTF8]{ctex}
\usepackage[a4paper,left=15mm,right=15mm,top=16mm,bottom=14mm,headheight=14pt,headsep=6pt,footskip=16pt]{geometry}
\usepackage{fancyhdr}
\pagestyle{fancy}
\fancyhf{}
\fancyhead[L]{\footnotesize ChabotAI伙伴 V1.4.3\quad 源程序}
\fancyhead[R]{\footnotesize 第 \thepage\ 页 / 共 60 页}
\fancyfoot[C]{\footnotesize 著作权人：陈聪}
\renewcommand{\headrulewidth}{0.4pt}
\renewcommand{\footrulewidth}{0.4pt}
\setlength{\parindent}{0pt}
\setlength{\parskip}{0pt}
\raggedbottom
\begin{document}

{\ttfamily\footnotesize\raggedright
//~================~FILE:~utils/chat-compress.js~================\\
/**\\
~*~聊天服务~----~压缩域：\\
~*~把上文交给~LLM~压缩为概要（手动触发~/~每轮回复后自动检查）\\
~*/\\
\mbox{}\\
import~\{~getChatRows,~getConversationCompression,~setConversationCompression~\}~from~'./storage.js'\\
import~\{~chatCompletion~\}~from~'./llm.js'\\
import~\{~addLog~\}~from~'./log.js'\\
import~\{~getConversationSettings~\}~from~'./chat-settings.js'\\
\mbox{}\\
const~COMPRESS\_KEEP\_TAIL~=~10~//~压缩时保留的最近完整消息条数\\
const~COMPRESS\_MIN\_NEW~=~4~//~每次压缩至少新增这么多条才执行\\
const~COMPRESS\_CHUNK~=~80~//~单次压缩请求处理的最大消息条数（超出分批多次调用）\\
const~COMPRESS\_SYSTEM~=~`你是对话压缩助手。把对话压缩成一份精炼概要，用于替代原文作为后续上下文。\\
要求：保留关键事实、用户偏好、已做的决定、约定和正在进行的事；保留重要人物/时间/地点信息；\\
用中文，150~字以内，只输出概要本身，不要解释、不要分段标题。`\\
\mbox{}\\
let~\_compressing~=~false\\
\mbox{}\\
/**\\
~*~把上文交给~LLM~压缩为概要（手动或自动触发）。\\
~*~压缩范围：上次压缩进度之后、保留尾部之前的消息；超过单批上限时分批调用。\\
~*~@param~\{boolean\}~force~手动触发传~true，忽略间隔/新增量限制\\
~*~@returns~\{Promise<boolean>\}~是否实际执行了压缩\\
~*/\\
export~async~function~compressContext(force~=~false)~\{\\
~~~~if~(\_compressing)~return~false\\
~~~~const~rows~=~getChatRows()\\
~~~~const~keepTail~=~COMPRESS\_KEEP\_TAIL\\
~~~~if~(rows.length~<=~keepTail~+~COMPRESS\_MIN\_NEW)~return~false\\
\mbox{}\\
~~~~const~\{~summary:~prev,~compressedUntil~\}~=~getConversationCompression()\\
~~~~const~since~=~Math.min(compressedUntil,~rows.length~-~keepTail)\\
~~~~const~cut~=~Math.max(since,~rows.length~-~keepTail)\\
~~~~if~(!force~\&\&~cut~-~since~<~COMPRESS\_MIN\_NEW)~return~false\\
~~~~if~(cut~<=~since)~return~false\\
\mbox{}\\
~~~~const~s~=~getConversationSettings()\\
~~~~if~(!s.apiKey~||~!s.baseUrl~||~!s.model)~return~false\\
\mbox{}\\
~~~~\_compressing~=~true\\
~~~~try~\{\\
~~~~~~~~let~merged~=~prev\\
~~~~~~~~let~cursor~=~since\\
~~~~~~~~//~分批压缩，后一批把前一批的概要合并进来，避免单次请求过大\\
~~~~~~~~while~(cursor~<~cut)~\{\\
~~~~~~~~~~~~const~end~=~Math.min(cursor~+~COMPRESS\_CHUNK,~cut)\\
~~~~~~~~~~~~const~batch~=~rows.slice(cursor,~end)\\
~~~~~~~~~~~~const~text~=~[
}

\newpage

{\ttfamily\footnotesize\raggedright
~~~~~~~~~~~~~~~~merged~?~`已有的对话概要：\textbackslash{}n\$\{merged\}`~:~'',\\
~~~~~~~~~~~~~~~~'以下是对话记录（越靠后越新）：',\\
~~~~~~~~~~~~~~~~batch.map((r)~=>~`\$\{r.role~===~'user'~?~'用户'~:~'你'\}：\$\{r.content\}`).join('\textbackslash{}n'),\\
~~~~~~~~~~~~~~~~'请把已有概要与新内容合并，输出一份完整的精炼概要。'\\
~~~~~~~~~~~~].filter(Boolean).join('\textbackslash{}n\textbackslash{}n')\\
~~~~~~~~~~~~const~resp~=~await~chatCompletion(\{\\
~~~~~~~~~~~~~~~~baseUrl:~s.baseUrl,\\
~~~~~~~~~~~~~~~~apiKey:~s.apiKey,\\
~~~~~~~~~~~~~~~~model:~s.model,\\
~~~~~~~~~~~~~~~~messages:~[\\
~~~~~~~~~~~~~~~~~~~~\{~role:~'system',~content:~COMPRESS\_SYSTEM~\},\\
~~~~~~~~~~~~~~~~~~~~\{~role:~'user',~content:~text~\}\\
~~~~~~~~~~~~~~~~],\\
~~~~~~~~~~~~~~~~temperature:~0.3,\\
~~~~~~~~~~~~~~~~maxTokens:~512,\\
~~~~~~~~~~~~~~~~reasoningEffort:~'none'~//~压缩任务简单，固定关闭思考，避免被思考~token~占满\\
~~~~~~~~~~~~\})\\
~~~~~~~~~~~~const~out~=~(resp.text~||~'').trim()\\
~~~~~~~~~~~~if~(!out)~return~false\\
~~~~~~~~~~~~merged~=~out\\
~~~~~~~~~~~~cursor~=~end\\
~~~~~~~~\}\\
~~~~~~~~setConversationCompression(merged,~cut)\\
~~~~~~~~addLog(\\
~~~~~~~~~~~~'info',\\
~~~~~~~~~~~~'压缩上文',\\
~~~~~~~~~~~~`压缩~\$\{cut~-~since\}~条消息为概要（保留尾部~\$\{keepTail\}~条）\textbackslash{}n---~压缩后上文~---\textbackslash{}n\$\{merg\\
ed\}`\\
~~~~~~~~)\\
~~~~~~~~return~true\\
~~~~\}~catch~(e)~\{\\
~~~~~~~~addLog('err',~'压缩上文失败',~e.message)\\
~~~~~~~~throw~e\\
~~~~\}~finally~\{\\
~~~~~~~~\_compressing~=~false\\
~~~~\}\\
\}\\
\mbox{}\\
/**\\
~*~自动压缩检查：距上次压缩新增消息达到间隔阈值则压缩。\\
~*~在每次对话回复落库后调用（异步、失败静默，不阻塞回复）。\\
~*/\\
export~async~function~maybeCompress()~\{\\
~~~~if~(\_compressing)~return~false\\
~~~~const~interval~=~getConversationSettings().compressInterval\\
~~~~if~(!interval)~return~false\\
~~~~const~\{~compressedUntil~\}~=~getConversationCompression()\\
~~~~if~(getChatRows().length~-~compressedUntil~<~interval)~return~false\\
~~~~addLog('info',~'自动压缩触发',~`间隔~\$\{interval\}~条`)\\
~~~~return~compressContext()
}

\newpage

{\ttfamily\footnotesize\raggedright
\}\\
\mbox{}\\
\mbox{}\\
//~================~FILE:~utils/chat-conversations.js~================\\
/**\\
~*~聊天服务~----~会话域：\\
~*~多会话管理（新建/切换/删除/复制）+~当前会话历史展示/清空~+~最近请求缓存同步\\
~*/\\
\mbox{}\\
import~\{\\
~~~~getChatRows,\\
~~~~clearChat,\\
~~~~getConversations,\\
~~~~getActiveConversationId,\\
~~~~createConversation,\\
~~~~switchConversation,\\
~~~~deleteConversation,\\
~~~~duplicateConversationToNew,\\
~~~~duplicateMemoriesToNew,\\
~~~~setConversationSettingsRaw\\
\}~from~'./storage.js'\\
import~\{~memoryStore,~lastRequest~\}~from~'./chat-state.js'\\
import~\{~getConversationSettings,~normalizeSettings~\}~from~'./chat-settings.js'\\
import~\{~addLog~\}~from~'./log.js'\\
\mbox{}\\
/**~剔除消息文本行首的~Scene/Memory~标记行（修复旧数据：早期落库过含标记的完整回复）~*/\\
function~stripMetaLines(content)~\{\\
~~~~const~clean~=~String(content~||~'')\\
~~~~~~~~.split('\textbackslash{}n')\\
~~~~~~~~.filter((l)~=>~!/\textasciicircum{}\textbackslash{}s*(Scene|Memory)\textbackslash{}s*[:：]/i.test(l))\\
~~~~~~~~.join('\textbackslash{}n')\\
~~~~~~~~.trim()\\
~~~~return~clean~||~'...'\\
\}\\
\mbox{}\\
/**~聊天页历史（返回副本，剔除遗留的~Scene/Memory~标记行后展示）~*/\\
export~function~getHistoryForUI()~\{\\
~~~~return~getChatRows().map((r)~=>~(\{~...r,~content:~stripMetaLines(r.content)~\}))\\
\}\\
\mbox{}\\
/**~清空对话~*/\\
export~function~clearConversation()~\{\\
~~~~clearChat()\\
~~~~syncLastRequest()\\
~~~~addLog('info',~'清空当前对话')\\
\}\\
\mbox{}\\
/**~同步请求缓存到当前会话的最后一条用户消息（切会话/清空后调用，避免跨会话重发错内容）~*/\\
export~function~syncLastRequest()~\{\\
~~~~const~rows~=~getChatRows()
}

\newpage

{\ttfamily\footnotesize\raggedright
~~~~for~(let~i~=~rows.length~-~1;~i~>=~0;~i--)~\{\\
~~~~~~~~if~(rows[i].role~===~'user')~\{\\
~~~~~~~~~~~~lastRequest.value~=~rows[i].content\\
~~~~~~~~~~~~return\\
~~~~~~~~\}\\
~~~~\}\\
~~~~lastRequest.value~=~''\\
\}\\
\mbox{}\\
//~----------~会话管理~----------\\
\mbox{}\\
/**~会话列表（供历史对话弹窗展示）~*/\\
export~function~listConversations()~\{\\
~~~~return~getConversations()\\
\}\\
\mbox{}\\
/**~当前会话~id~*/\\
export~function~activeConversationId()~\{\\
~~~~return~getActiveConversationId()\\
\}\\
\mbox{}\\
/**~开始新对话：当前对话非空则归档并新建，否则复用并重置当前空对话；新对话复制当前设置（每个对话独享\\
一份设置）~*/\\
export~function~startNewConversation()~\{\\
~~~~memoryStore.resetState()\\
~~~~//~记录创建前的当前会话生效设置（含无快照会话的全局回退值），供新对话复制\\
~~~~const~currentSettings~=~getConversationSettings()\\
~~~~if~(getChatRows().length)~\{\\
~~~~~~~~createConversation()\\
~~~~~~~~addLog('info',~'开始新对话')\\
~~~~\}~else~\{\\
~~~~~~~~clearChat()\\
~~~~~~~~addLog('info',~'重置空对话')\\
~~~~\}\\
~~~~setConversationSettingsRaw(normalizeSettings(currentSettings))\\
~~~~syncLastRequest()\\
\}\\
\mbox{}\\
/**~切换到指定会话（记忆与人格随会话切换，重置召回状态）~*/\\
export~function~openConversation(id)~\{\\
~~~~const~ok~=~switchConversation(id)\\
~~~~if~(ok)~\{\\
~~~~~~~~memoryStore.resetState()\\
~~~~~~~~syncLastRequest()\\
~~~~~~~~addLog('info',~'切换会话',~id)\\
~~~~\}\\
~~~~return~ok\\
\}\\
\mbox{}\\
/**~删除指定会话~*/
}

\newpage

{\ttfamily\footnotesize\raggedright
export~function~removeConversation(id)~\{\\
~~~~const~ok~=~deleteConversation(id)\\
~~~~if~(ok)~\{\\
~~~~~~~~memoryStore.resetState()\\
~~~~~~~~syncLastRequest()\\
~~~~~~~~addLog('info',~'删除会话',~id)\\
~~~~\}\\
~~~~return~ok\\
\}\\
\mbox{}\\
/**~复制当前会话（消息+记忆+人格）到新会话并切换~*/\\
export~function~copyConversationToNew()~\{\\
~~~~duplicateConversationToNew()\\
~~~~memoryStore.resetState()\\
~~~~syncLastRequest()\\
~~~~addLog('info',~'复制对话到新会话')\\
\}\\
\mbox{}\\
/**~复制当前会话的记忆到新会话并切换~*/\\
export~function~copyMemoriesToNew()~\{\\
~~~~duplicateMemoriesToNew()\\
~~~~memoryStore.resetState()\\
~~~~syncLastRequest()\\
~~~~addLog('info',~'复制记忆到新会话')\\
\}\\
\mbox{}\\
\mbox{}\\
//~================~FILE:~utils/chat-proactive.js~================\\
/**\\
~*~聊天服务~----~拟真聊天域：\\
~*~前台定时器调度（随机时间节点主动发消息）+~回前台补发（catch-up）+~主动消息发送链路。\\
~*\\
~*~生效条件（当前会话）：proactiveEnabled~开启~且~timeMode=real（现实时间）。\\
~*~低后台占用设计：单条链式~setTimeout（非轮询，空闲零唤醒），退后台**不清定时器**----\\
~*~单条~pending~timeout~后台零~CPU；H5~隐藏标签页会被浏览器节流（约~1~次/分钟，够用）；\\
~*~App~后台~JS~挂起时定时器不触发，回前台由~catchUpProactive~检测到期补发。\\
~*~**进程结束后重开（关页/杀进程）也不丢触发**：到期时刻持久化到~storage，\\
~*~重开时~catchUpProactive~据此判断"关闭期间已到触发时刻"，已过则补发一次再重排。\\
~*\\
~*~主动消息格式（与普通回复一致，由~parseAndValidateReply~强制校验）：\\
~*~每条~<=1~句、连续表情~<=2，多条消息用换行分隔；话题临近结束时模型用短句/表情收尾。\\
~*/\\
\mbox{}\\
import~\{\\
~~~~getChatRows,\\
~~~~addChatRow,\\
~~~~getSceneHistory,\\
~~~~setScene,\\
~~~~getConversationCompression,\\
~~~~getSetting,
}

\newpage

{\ttfamily\footnotesize\raggedright
~~~~setSetting\\
\}~from~'./storage.js'\\
import~\{~memoryStore~\}~from~'./chat-state.js'\\
import~\{~getConversationSettings~\}~from~'./chat-settings.js'\\
import~\{~buildSystemPrompt,~buildNowText,~getPersonalityById~\}~from~'./prompts.js'\\
import~\{~chatCompletion~\}~from~'./llm.js'\\
import~\{~addLog~\}~from~'./log.js'\\
import~\{~emojiListForPrompt~\}~from~'./emojis.js'\\
import~\{~maybeCompress~\}~from~'./chat-compress.js'\\
import~\{~parseAndValidateReply~\}~from~'./chat.js'\\
import~\{~notifyProactive~\}~from~'./notify.js'\\
\mbox{}\\
//~----------~常量~----------\\
\mbox{}\\
/**~频率档位~->~随机间隔区间（分钟）~*/\\
const~LEVEL\_RANGES~=~\{\\
~~~~low:~[45,~120],\\
~~~~medium:~[15,~45],\\
~~~~high:~[5,~15]\\
\}\\
const~TICK\_MAX\_MS~=~60~*~60~*~1000~//~单次定时最长~60~分钟：即使间隔更长也先醒一次，保证窗口/设置变\\
化及时响应\\
const~HISTORY\_ENTRIES~=~15~//~主动消息注入的对话历史条数（上限，与~sendMessage~一致）\\
const~PROACTIVE\_DUE\_KEY~=~'proactive\_due\_at'~//~到期时刻持久化键：进程结束（关页/杀进程）后重开时，\\
据此感知"关闭期间已到触发时刻"并补发\\
\mbox{}\\
//~----------~状态~----------\\
\mbox{}\\
let~\_timer~=~null~//~当前定时器\\
let~\_dueAt~=~0~//~下一触发时刻（退后台期间保留，供回前台补发判断）\\
let~\_suppressed~=~false~//~聊天页~loading~时抑制（等待~LLM~回复期间不打扰）\\
let~\_busy~=~false~//~主动请求进行中（防止重入）\\
\mbox{}\\
/**~读取持久化的到期时刻（绝对时间戳，无则~0）~*/\\
function~\_loadDueAt()~\{\\
~~~~return~parseInt(getSetting(PROACTIVE\_DUE\_KEY,~0),~10)~||~0\\
\}\\
\mbox{}\\
/**~持久化到期时刻（进程结束后重开时用于判断是否错过触发；0~=~清空）~*/\\
function~\_saveDueAt(v)~\{\\
~~~~try~\{\\
~~~~~~~~setSetting(PROACTIVE\_DUE\_KEY,~v)\\
~~~~\}~catch~(e)~\{\\
~~~~~~~~/*~ignore：存储失败不影响调度本身~*/\\
~~~~\}\\
\}\\
\mbox{}\\
function~\_clearTimer()~\{\\
~~~~if~(\_timer)~\{\\
~~~~~~~~clearTimeout(\_timer)
}

\newpage

{\ttfamily\footnotesize\raggedright
~~~~~~~~\_timer~=~null\\
~~~~\}\\
\}\\
\mbox{}\\
/**~聊天页同步~loading~状态：true~时调度器到期不发（用户在等~LLM~回复）~*/\\
export~function~setProactiveSuppressed(v)~\{\\
~~~~const~was~=~\_suppressed\\
~~~~\_suppressed~=~!!v\\
~~~~//~解除抑制时若调度器已失活（如~loading~期间到期未重排），立即重排，避免此后不再触发\\
~~~~if~(!\_suppressed~\&\&~was~\&\&~!\_timer)~\_schedule()\\
\}\\
\mbox{}\\
/**\\
~*~功能是否开启（与瞬时抑制无关）：开关~+~现实时间~+~API~配置齐备。\\
~*~调度重排只看此项----只要功能开着就持续重排，loading~抑制只拦截"到期发送"不拦截重排，\\
~*~保证任何时序下调度器都不会失活（否则会表现为"等待调度"、只能重启恢复）。\\
~*/\\
function~\_featureOn(s)~\{\\
~~~~return~!!(\\
~~~~~~~~s~\&\&\\
~~~~~~~~s.proactiveEnabled~===~true~\&\&\\
~~~~~~~~s.timeMode~===~'real'~\&\&\\
~~~~~~~~s.apiKey~\&\&\\
~~~~~~~~s.baseUrl~\&\&\\
~~~~~~~~s.model\\
~~~~)\\
\}\\
\mbox{}\\
/**~当前会话是否具备主动发送的基本条件（功能开启~+~未被~loading~抑制）~*/\\
function~\_canRun(s)~\{\\
~~~~return~\_featureOn(s)~\&\&~!\_suppressed\\
\}\\
\mbox{}\\
/**~最近一条用户消息时间（无则~0）~*/\\
function~\_lastUserAt()~\{\\
~~~~const~rows~=~getChatRows()\\
~~~~for~(let~i~=~rows.length~-~1;~i~>=~0;~i--)~\{\\
~~~~~~~~if~(rows[i].role~===~'user')~return~Date.parse(rows[i].created\_at)~||~0\\
~~~~\}\\
~~~~return~0\\
\}\\
\mbox{}\\
/**~当前时刻是否在主动消息时段窗口内（分钟级，闭区间~[start,~end]，start/end~为当天第几分钟~0-1439）\\
~*/\\
function~\_inWindow(s)~\{\\
~~~~const~d~=~new~Date()\\
~~~~const~min~=~d.getHours()~*~60~+~d.getMinutes()\\
~~~~return~min~>=~s.proactiveStartMin~\&\&~min~<=~s.proactiveEndMin\\
\}\\
\mbox{}
}

\newpage

{\ttfamily\footnotesize\raggedright
/**~开关自定义倒计时（调试值~>0~即处于"测试模式"）：调试模式下忽略时段窗口与"必须有用户消息"门禁，便\\
于随时验证链路~*/\\
function~\_isDebugCountdown(s)~\{\\
~~~~return~parseInt(s~\&\&~s.proactiveCustomSeconds,~10)~>~0\\
\}\\
\mbox{}\\
/**~主动消息被门禁拦截时写一条可见日志，避免"只见调度、不知为何不发"~*/\\
function~\_skipLog(reason)~\{\\
~~~~addLog('info',~'主动消息跳过',~reason)\\
\}\\
\mbox{}\\
/**~当天第几分钟~->~"HH:MM"~展示~*/\\
function~\_fmtClockMin(min)~\{\\
~~~~const~h~=~String(Math.floor(min~/~60)).padStart(2,~'0')\\
~~~~const~m~=~String(min~\%~60).padStart(2,~'0')\\
~~~~return~`\$\{h\}:\$\{m\}`\\
\}\\
\mbox{}\\
/**~下次触发延迟（ms）：调试用自定义倒计时（proactiveCustomSeconds>0）优先；否则窗口外等到窗口起点、\\
窗口内按档位随机间隔。始终~>=10s（自定义）或~>=60s，防忙循环/防误触~API~轰炸~*/\\
function~\_nextDelay(s)~\{\\
~~~~const~custom~=~parseInt(s~\&\&~s.proactiveCustomSeconds,~10)~||~0\\
~~~~if~(custom~>~0)~return~Math.max(10000,~custom~*~1000)\\
~~~~const~d~=~new~Date()\\
~~~~const~min~=~d.getHours()~*~60~+~d.getMinutes()\\
~~~~const~[lo,~hi]~=~LEVEL\_RANGES[s.proactiveLevel]~||~LEVEL\_RANGES.medium\\
~~~~const~toStart~=~(targetMin)~=>~\{\\
~~~~~~~~//~从当前时刻算起到~targetMin~这一分钟开头的毫秒数（同分钟时为负，代表"已过该点"）\\
~~~~~~~~return~(targetMin~-~min)~*~60000~-~d.getSeconds()~*~1000~-~d.getMilliseconds()\\
~~~~\}\\
~~~~if~(min~<~s.proactiveStartMin)~\{\\
~~~~~~~~return~Math.max(60000,~toStart(s.proactiveStartMin))\\
~~~~\}\\
~~~~if~(min~>~s.proactiveEndMin)~\{\\
~~~~~~~~return~Math.max(60000,~toStart(s.proactiveStartMin~+~1440))\\
~~~~\}\\
~~~~return~Math.max(60000,~(lo~+~Math.random()~*~(hi~-~lo))~*~60000)\\
\}\\
\mbox{}\\
//~----------~调度器~----------\\
\mbox{}\\
/**~毫秒~->~"X~小时~Y~分"~/~"X~分~Y~秒"~/~"X~秒"（调度日志与倒计时展示用）~*/\\
function~\_fmtDuration(ms)~\{\\
~~~~const~total~=~Math.max(0,~Math.ceil((ms~||~0)~/~1000))\\
~~~~const~h~=~Math.floor(total~/~3600)\\
~~~~const~m~=~Math.floor((total~\%~3600)~/~60)\\
~~~~const~s~=~total~\%~60\\
~~~~if~(h~>~0)~return~`\$\{h\}~小时~\$\{m\}~分`\\
~~~~if~(m~>~0)~return~`\$\{m\}~分~\$\{String(s).padStart(2,~'0')\}~秒`\\
~~~~return~`\$\{s\}~秒`
}

\newpage

{\ttfamily\footnotesize\raggedright
\}\\
\mbox{}\\
/**~按当前会话设置安排下一次触发（清旧定时器后重排），并在调试日志记录调度状态。只看功能开关，抑制/\\
时段只拦发送不拦重排~*/\\
function~\_schedule()~\{\\
~~~~\_clearTimer()\\
~~~~const~s~=~getConversationSettings()\\
~~~~if~(!\_featureOn(s))~\{\\
~~~~~~~~//~功能关闭（关闭拟真/切虚拟时间/未配~API）：清掉持久化的到期时刻，避免残留值在日后重开时误\\
触发\\
~~~~~~~~\_dueAt~=~0\\
~~~~~~~~\_saveDueAt(0)\\
~~~~~~~~return\\
~~~~\}\\
~~~~const~delay~=~\_nextDelay(s)\\
~~~~\_dueAt~=~Date.now()~+~delay\\
~~~~\_saveDueAt(\_dueAt)\\
~~~~\_timer~=~setTimeout(\_tick,~Math.min(delay,~TICK\_MAX\_MS))\\
~~~~addLog('info',~'拟真聊天调度',~`下次主动消息：\$\{\_fmtDuration(delay)\}~后（\$\{s.proactiveLevel\}~档\\
）`)\\
\}\\
\mbox{}\\
/**~进程重启后恢复定时器：持久化到期时刻仍在未来->按剩余时间重挂；否则全新重排（或已补发过->重新起一\\
轮）~*/\\
function~\_armFromPersisted()~\{\\
~~~~if~(\_dueAt~\&\&~\_dueAt~>~Date.now())~\{\\
~~~~~~~~\_timer~=~setTimeout(\_tick,~Math.min(\_dueAt~-~Date.now(),~TICK\_MAX\_MS))\\
~~~~\}~else~\{\\
~~~~~~~~\_schedule()\\
~~~~\}\\
\}\\
\mbox{}\\
function~\_tick()~\{\\
~~~~\_timer~=~null\\
~~~~//~\_dueAt~为~0~说明已由~catchUp~补发过，此处不应再次触发（否则重复发送）\\
~~~~if~(\_dueAt~\&\&~Date.now()~>=~\_dueAt)~\{\\
~~~~~~~~\_dueAt~=~0\\
~~~~~~~~//~到期触发：每次触发都必定产生一条可见日志----\\
~~~~~~~~//~~~成功~->~主动消息已发送（+~LLM~请求/响应）；\\
~~~~~~~~//~~~被门禁拦截~->~主动消息跳过（附原因），以及「拟真聊天到期.未发送」标记；\\
~~~~~~~~//~~~异常~->~主动消息失败（附堆栈），杜绝"只见调度、无任何输出"的静默失败。\\
~~~~~~~~sendProactiveBurst()\\
~~~~~~~~~~~~.then((r)~=>~\{\\
~~~~~~~~~~~~~~~~if~(r~===~null)~addLog('info',~'拟真聊天到期',~'已到触发时刻但未发送（原因见「主动消\\
息跳过」日志）')\\
~~~~~~~~~~~~\})\\
~~~~~~~~~~~~.catch((e)~=>~addLog('err',~'主动消息失败',~(e~\&\&~e.stack)~||~String(e)))\\
~~~~\}\\
~~~~\_schedule()\\
\}
}

\newpage

{\ttfamily\footnotesize\raggedright
\mbox{}\\
/**\\
~*~下次主动消息倒计时（毫秒）；未调度（功能未开启/尚未重排）返回~null。\\
~*~供设置页调试区实时显示调度状态。\\
~*/\\
export~function~getProactiveCountdown()~\{\\
~~~~if~(!\_timer~||~!\_dueAt)~return~null\\
~~~~const~s~=~getConversationSettings()\\
~~~~if~(!\_featureOn(s))~return~null\\
~~~~return~Math.max(0,~\_dueAt~-~Date.now())\\
\}\\
\mbox{}\\
/**\\
~*~回前台/重开补发：若到期时刻已过（含进程结束后重开的场景），立即补发一次；随后恢复调度。\\
~*~幂等，可随时调用（切会话/改设置后用于让调度器立即按新状态重排）。\\
~*/\\
export~function~catchUpProactive()~\{\\
~~~~const~s~=~getConversationSettings()\\
~~~~if~(!\_featureOn(s))~\{\\
~~~~~~~~//~功能关闭（关闭拟真/切到虚拟时间/未配~API）：清掉在途定时器与持久化到期时刻\\
~~~~~~~~\_clearTimer()\\
~~~~~~~~\_dueAt~=~0\\
~~~~~~~~\_saveDueAt(0)\\
~~~~~~~~return\\
~~~~\}\\
~~~~//~进程结束后内存~\_dueAt~清零，但持久化的到期时刻仍在：据其判断"关闭期间是否已到触发时刻"\\
~~~~if~(!\_dueAt)~\_dueAt~=~\_loadDueAt()\\
~~~~if~(\_dueAt~\&\&~Date.now()~>=~\_dueAt)~\{\\
~~~~~~~~//~清掉即将到期的旧定时器，避免其随后再触发一次造成重复发送，并清空持久化到期时刻\\
~~~~~~~~\_clearTimer()\\
~~~~~~~~\_dueAt~=~0\\
~~~~~~~~\_saveDueAt(0)\\
~~~~~~~~sendProactiveBurst()\\
~~~~~~~~~~~~.then((r)~=>~\{\\
~~~~~~~~~~~~~~~~if~(r~===~null)~addLog('info',~'拟真聊天补发',~'已到触发时刻但未发送（原因见「主动消\\
息跳过」日志）')\\
~~~~~~~~~~~~\})\\
~~~~~~~~~~~~.catch((e)~=>~addLog('err',~'主动消息失败',~(e~\&\&~e.stack)~||~String(e)))\\
~~~~\}\\
~~~~if~(!\_timer)~\_armFromPersisted()\\
\}\\
\mbox{}\\
/**\\
~*~立即按当前会话设置重排下一次触发（修改倒计时/档位/时段/切会话后调用）。\\
~*~与~catchUpProactive~不同：即使定时器已挂载也强制重算（catchUp~为保持跨页倒计时不重置而跳过）。\\
~*/\\
export~function~rearmProactive()~\{\\
~~~~\_schedule()\\
\}\\
\mbox{}
}

\newpage

{\ttfamily\footnotesize\raggedright
//~----------~主动消息发送~----------\\
\mbox{}\\
/**\\
~*~发送一条主动消息（拟真~burst）：组装~system（含拟真规则）+~最近历史，调~LLM，\\
~*~校验后按换行拆多条落库（Scene/Memory~由校验写入），通知聊天页刷新。\\
~*~@param~\{\{force?:~boolean\}\}~[opts]~force=true~用于调试按钮：忽略时段/空闲/间隔门禁，仅校验~API~配\\
置\\
~*~@returns~\{Promise<\{lines:string[],~saved:number\}|null>\}~未满足触发条件返回~null\\
~*/\\
export~async~function~sendProactiveBurst(opts~=~\{\})~\{\\
~~~~const~force~=~!!opts.force\\
~~~~const~s~=~getConversationSettings()\\
~~~~const~debugCountdown~=~\_isDebugCountdown(s)\\
~~~~//~基本条件（开关~+~现实时间~+~未抑制）仅非调试时要求；调试按钮忽略开关/时段/空闲/抑制，但仍要求\\
~API~配置\\
~~~~if~(!force~\&\&~!\_canRun(s))~\{\\
~~~~~~~~\_skipLog(\_suppressed~?~'聊天页正在等待~LLM~回复（已抑制），本轮跳过'~:~'功能未开启（需开启拟\\
真~+~现实时间~+~已配置~API）')\\
~~~~~~~~return~null\\
~~~~\}\\
~~~~if~(!s.apiKey~||~!s.baseUrl~||~!s.model)~\{\\
~~~~~~~~\_skipLog('API~配置不完整（缺少~baseUrl/API~Key/model）')\\
~~~~~~~~return~null\\
~~~~\}\\
~~~~//~触发门禁：时段窗口~+~已有用户消息。自定义倒计时为调试工具，忽略这两项门禁，以便随时测试\\
~~~~if~(!force)~\{\\
~~~~~~~~if~(!debugCountdown~\&\&~!\_inWindow(s))~\{\\
~~~~~~~~~~~~\_skipLog(`不在主动消息时段窗口（\$\{\_fmtClockMin(s.proactiveStartMin)\}-\$\{\_fmtClockMin(s.pr\\
oactiveEndMin)\}）`)\\
~~~~~~~~~~~~return~null\\
~~~~~~~~\}\\
~~~~~~~~if~(!debugCountdown~\&\&~!\_lastUserAt())~\{\\
~~~~~~~~~~~~\_skipLog('当前会话还没有用户消息，本轮不主动发送')\\
~~~~~~~~~~~~return~null\\
~~~~~~~~\}\\
~~~~\}\\
~~~~if~(\_busy)~\{\\
~~~~~~~~\_skipLog('上一条主动消息仍在处理中，本轮跳过')\\
~~~~~~~~return~null\\
~~~~\}\\
~~~~\_busy~=~true\\
~~~~try~\{\\
~~~~~~~~//~组装~system（与~sendMessage~一致，proactive~注入拟真规则）\\
~~~~~~~~const~memoryText~=~memoryStore.retrieveContext('')\\
~~~~~~~~const~personalityPrompt~=\\
~~~~~~~~~~~~s.personalityId~===~'custom'~?~s.customPrompt~:~getPersonalityById(s.personalityId).prom\\
pt\\
~~~~~~~~const~system~=~buildSystemPrompt(\\
~~~~~~~~~~~~personalityPrompt~||~'你是友好的聊天伙伴。',\\
~~~~~~~~~~~~memoryText,
}

\newpage

{\ttfamily\footnotesize\raggedright
~~~~~~~~~~~~getSceneHistory(),\\
~~~~~~~~~~~~s.timeMode~===~'virtual'~?~''~:~buildNowText(),\\
~~~~~~~~~~~~s.emojiEnabled~?~emojiListForPrompt()~:~[],\\
~~~~~~~~~~~~memoryStore.l1Usage(),\\
~~~~~~~~~~~~true\\
~~~~~~~~)\\
~~~~~~~~const~\{~summary,~compressedUntil~\}~=~getConversationCompression()\\
~~~~~~~~const~messages~=~[\\
~~~~~~~~~~~~\{~role:~'system',~content:~summary~?~`\$\{system\}\textbackslash{}n\textbackslash{}n此前对话概要：\$\{summary\}`~:~system~\}\\
~~~~~~~~]\\
~~~~~~~~for~(const~h~of~getChatRows().slice(compressedUntil).slice(-HISTORY\_ENTRIES))~\{\\
~~~~~~~~~~~~messages.push(\{~role:~h.role,~content:~h.content~\})\\
~~~~~~~~\}\\
~~~~~~~~//~末尾追加内部指令（仅请求，不落库）：让模型以"主动开口"的方式输出下一条消息\\
~~~~~~~~messages.push(\{~role:~'user',~content:~'（内部指令：请作为~AI~主动开口给用户发消息，只输出这\\
条主动消息本身）'~\})\\
\mbox{}\\
~~~~~~~~//~请求~->~拟真格式校验（每条~<=1~句、连续表情~<=2）->~不合格自动重试\\
~~~~~~~~const~maxAttempts~=~Math.max(1,~Math.min(20,~parseInt(s.maxRequestAttempts,~10)~||~5))\\
~~~~~~~~let~result~=~null\\
~~~~~~~~let~lastReason~=~''\\
~~~~~~~~for~(let~attempt~=~1;~attempt~<=~maxAttempts;~attempt++)~\{\\
~~~~~~~~~~~~const~reply~=~await~chatCompletion(\{\\
~~~~~~~~~~~~~~~~baseUrl:~s.baseUrl,\\
~~~~~~~~~~~~~~~~apiKey:~s.apiKey,\\
~~~~~~~~~~~~~~~~model:~s.model,\\
~~~~~~~~~~~~~~~~messages,\\
~~~~~~~~~~~~~~~~temperature:~s.temperature,\\
~~~~~~~~~~~~~~~~reasoningEffort:~s.reasoningEffort\\
~~~~~~~~~~~~\})\\
~~~~~~~~~~~~const~parsed~=~parseAndValidateReply(reply.text,~\{~proactive:~true~\})\\
~~~~~~~~~~~~if~(parsed.ok)~\{\\
~~~~~~~~~~~~~~~~result~=~parsed\\
~~~~~~~~~~~~~~~~break\\
~~~~~~~~~~~~\}\\
~~~~~~~~~~~~lastReason~=~parsed.reason\\
~~~~~~~~~~~~addLog('info',~'主动消息格式不合格，自动重试',~`第~\$\{attempt\}/\$\{maxAttempts\}~次：\$\{lastR\\
eason\}`)\\
~~~~~~~~\}\\
~~~~~~~~if~(!result)~\{\\
~~~~~~~~~~~~addLog('err',~'主动消息格式不合格（已达请求上限）',~`最大请求次数~\$\{maxAttempts\}：\$\{last\\
Reason\}`)\\
~~~~~~~~~~~~throw~new~Error(`主动消息格式连续~\$\{maxAttempts\}~次不合格（\$\{lastReason\}）`)\\
~~~~~~~~\}\\
\mbox{}\\
~~~~~~~~//~按换行拆分多条落库（每条~<=1~句一条气泡），rollback~挂最后一行\\
~~~~~~~~const~lines~=~(result.cleanReply~||~'').split('\textbackslash{}n').map((l)~=>~l.trim()).filter(Boolean)\\
~~~~~~~~if~(!lines.length)~lines.push('...')\\
~~~~~~~~const~sceneLenBefore~=~getSceneHistory().length\\
~~~~~~~~const~rollback~=~\{~sceneLenBefore,~memoryUndos:~result.memoryUndos~||~[]~\}
}

\newpage

{\ttfamily\footnotesize\raggedright
~~~~~~~~lines.forEach((line,~idx)~=>~\{\\
~~~~~~~~~~~~addChatRow('assistant',~line,~idx~===~lines.length~-~1~?~\{~rollback~\}~:~undefined)\\
~~~~~~~~\})\\
~~~~~~~~if~(result.newScene)~\{\\
~~~~~~~~~~~~setScene(result.newScene)\\
~~~~~~~~~~~~addLog('info',~'主动消息情景更新',~result.newScene)\\
~~~~~~~~\}\\
~~~~~~~~if~(result.saved~>~0)~addLog('info',~`主动消息记忆入库~\$\{result.saved\}~条`)\\
~~~~~~~~memoryStore.maintenance()\\
~~~~~~~~maybeCompress().catch(()~=>~\{~\})\\
~~~~~~~~addLog('info',~'主动消息已发送',~lines.join('~/~'))\\
~~~~~~~~//~系统通知：仅在应用处于后台时弹（前台刚发完用户正在看，不打扰）\\
~~~~~~~~notifyProactive(lines.join('\textbackslash{}n'))\\
~~~~~~~~//~通知聊天页刷新（uni.\$emit~在~App/H5/小程序均可用；Node~测试环境无此~API，跳过）\\
~~~~~~~~if~(typeof~uni~!==~'undefined'~\&\&~uni.\$emit)~uni.\$emit('proactive-burst',~\{~lines~\})\\
~~~~~~~~return~\{~lines,~saved:~result.saved~\}\\
~~~~\}~finally~\{\\
~~~~~~~~\_busy~=~false\\
~~~~\}\\
\}\\
\mbox{}\\
/**~调试按钮：立即发送一条主动消息（忽略时段/空闲/间隔门禁，仅校验~API~配置与开关）~*/\\
export~async~function~debugProactiveMessage()~\{\\
~~~~return~sendProactiveBurst(\{~force:~true~\})\\
\}\\
\mbox{}\\
\mbox{}\\
//~================~FILE:~utils/chat-settings.js~================\\
/**\\
~*~聊天服务~----~设置域：\\
~*~全局默认设置读取~/~会话设置快照读写（每个对话独享一份）/~会话人格保存\\
~*/\\
\mbox{}\\
import~\{\\
~~~~getSetting,\\
~~~~getConversationSettingsRaw,\\
~~~~setConversationSettingsRaw,\\
~~~~getConversationPersonality\\
\}~from~'./storage.js'\\
import~\{~getPersonaName~\}~from~'./prompts.js'\\
import~\{~addLog~\}~from~'./log.js'\\
\mbox{}\\
/**~人格名称（含自定义，供~UI~展示）~*/\\
export~function~personaName(id)~\{\\
~~~~return~getPersonaName(id)\\
\}\\
\mbox{}\\
/**~全局默认设置（无设置快照的会话回退用；设置面板编辑的是当前会话设置）~*/\\
export~function~getSettings()~\{\\
~~~~const~personalityId~=~getSetting('personalityId',~'gentle')
}

\newpage

{\ttfamily\footnotesize\raggedright
~~~~return~\{\\
~~~~~~~~baseUrl:~getSetting('baseUrl',~'https://api.openai.com/v1'),\\
~~~~~~~~apiKey:~getSetting('apiKey',~''),\\
~~~~~~~~model:~getSetting('model',~'gpt-5.4-mini'),\\
~~~~~~~~temperature:~parseFloat(getSetting('temperature',~'0.8')),\\
~~~~~~~~reasoningEffort:~getSetting('reasoningEffort',~'none'),~//~思考模式：none~关闭~/~high~开启~/\\
~''~跟随模型（Ollama~等兼容接口经~reasoning\_effort~控制）\\
~~~~~~~~personalityId,\\
~~~~~~~~customPrompt:~getSetting('customPrompt',~''),\\
~~~~~~~~timeMode:~getSetting('timeMode',~'real'),~//~情景时间模式：real~现实时间~/~virtual~虚拟时间\\
~~~~~~~~compressInterval:~parseInt(getSetting('compressInterval',~'0'),~10)~||~0,~//~自动压缩间隔（\\
条），0=关闭\\
~~~~~~~~maxRequestAttempts:~parseInt(getSetting('maxRequestAttempts',~'5'),~10)~||~5,~//~回复格式不\\
合格时的最大请求次数（重试上限）\\
~~~~~~~~emojiEnabled:~getSetting('emojiEnabled',~true)~!==~false,~//~聊天表情包开关：关闭后请求不携\\
带表情清单\\
~~~~~~~~//~语音阅读（TTS）：默认关闭；开启后~LLM~新回复会自动合成语音播放\\
~~~~~~~~ttsEnabled:~getSetting('ttsEnabled',~false)~!==~false,\\
~~~~~~~~ttsApiKey:~getSetting('ttsApiKey',~''),\\
~~~~~~~~ttsModel:~getSetting('ttsModel',~'qwen3-tts-flash'),\\
~~~~~~~~ttsVoice:~getSetting('ttsVoice',~'Cherry'),\\
~~~~~~~~//~拟真聊天：仅现实时间模式生效；AI~在随机时间主动发消息（每条~<=1~句、连续表情~<=2）\\
~~~~~~~~proactiveEnabled:~getSetting('proactiveEnabled',~false),\\
~~~~~~~~proactiveStartMin:~getSetting('proactiveStartMin',~9~*~60),~//~主动消息时段起（当天第几分钟~\\
0-1439，默认~09:00）\\
~~~~~~~~proactiveEndMin:~getSetting('proactiveEndMin',~23~*~60~+~59),~//~主动消息时段止（当天第几分\\
钟~0-1439，默认~23:59）\\
~~~~~~~~proactiveLevel:~getSetting('proactiveLevel',~'medium'),~//~频率档位~low/medium/high\\
~~~~~~~~proactiveCustomSeconds:~parseInt(getSetting('proactiveCustomSeconds',~'0'),~10)~||~0,~//~调\\
试：自定义倒计时（秒），>0~时覆盖档位\\
~~~~~~~~bubbleOpacity:~(()~=>~\{\\
~~~~~~~~~~~~const~n~=~parseFloat(getSetting('bubbleOpacity',~'1'))\\
~~~~~~~~~~~~return~Number.isFinite(n)~?~Math.max(0.2,~Math.min(1,~n))~:~1\\
~~~~~~~~\})(),~//~聊天气泡背景不透明度~0.2\textasciitilde{}1（1=不透明）\\
~~~~~~~~personaName:~getPersonaName(personalityId)\\
~~~~\}\\
\}\\
\mbox{}\\
/**~归一化设置对象（saveSettings~/~saveConversationPersonality~/~新对话复制统一调用）~*/\\
export~function~normalizeSettings(s)~\{\\
~~~~return~\{\\
~~~~~~~~baseUrl:~String(s~\&\&~s.baseUrl~?~s.baseUrl~:~'').trim(),\\
~~~~~~~~apiKey:~String(s~\&\&~s.apiKey~?~s.apiKey~:~'').trim(),\\
~~~~~~~~model:~String(s~\&\&~s.model~?~s.model~:~'').trim(),\\
~~~~~~~~temperature:~Number.isFinite(parseFloat(s~\&\&~s.temperature))~?~parseFloat(s.temperature)~:~0\\
.8,\\
~~~~~~~~reasoningEffort:~(s~\&\&~s.reasoningEffort)~||~'',\\
~~~~~~~~personalityId:~(s~\&\&~s.personalityId)~||~'gentle',\\
~~~~~~~~customPrompt:~String(s~\&\&~s.customPrompt~?~s.customPrompt~:~'').trim(),\\
~~~~~~~~timeMode:~s~\&\&~s.timeMode~===~'virtual'~?~'virtual'~:~'real',
}

\newpage

{\ttfamily\footnotesize\raggedright
~~~~~~~~compressInterval:~parseInt(s~\&\&~s.compressInterval,~10)~||~0,\\
~~~~~~~~maxRequestAttempts:~Math.max(1,~Math.min(20,~parseInt(s~\&\&~s.maxRequestAttempts,~10)~||~5)),\\
~~~~~~~~emojiEnabled:~!s~||~s.emojiEnabled~!==~false,\\
~~~~~~~~ttsEnabled:~!s~||~s.ttsEnabled~!==~false,\\
~~~~~~~~ttsApiKey:~String(s~\&\&~s.ttsApiKey~?~s.ttsApiKey~:~'').trim(),\\
~~~~~~~~ttsModel:~String(s~\&\&~s.ttsModel~?~s.ttsModel~:~'').trim()~||~'qwen3-tts-flash',\\
~~~~~~~~ttsVoice:~String(s~\&\&~s.ttsVoice~?~s.ttsVoice~:~'').trim()~||~'Cherry',\\
~~~~~~~~//~拟真聊天：开关默认关；时段为分钟级（当天第几分钟~0-1439，默认~09:00-23:59）；档位白名单\\
~~~~~~~~proactiveEnabled:~!!s~\&\&~s.proactiveEnabled~===~true,\\
~~~~~~~~proactiveStartMin:~(()~=>~\{\\
~~~~~~~~~~~~const~n~=~parseInt(s~\&\&~s.proactiveStartMin,~10)\\
~~~~~~~~~~~~return~Number.isNaN(n)~?~9~*~60~:~Math.max(0,~Math.min(1439,~n))\\
~~~~~~~~\})(),\\
~~~~~~~~proactiveEndMin:~(()~=>~\{\\
~~~~~~~~~~~~const~n~=~parseInt(s~\&\&~s.proactiveEndMin,~10)\\
~~~~~~~~~~~~return~Number.isNaN(n)~?~23~*~60~+~59~:~Math.max(0,~Math.min(1439,~n))\\
~~~~~~~~\})(),\\
~~~~~~~~proactiveLevel:~['low',~'medium',~'high'].includes(s~\&\&~s.proactiveLevel)~?~s.proactiveLevel\\
~:~'medium',\\
~~~~~~~~//~自定义倒计时（秒）：0=用档位；钳制在~0-3600，供调试快速触发\\
~~~~~~~~proactiveCustomSeconds:~(()~=>~\{\\
~~~~~~~~~~~~const~n~=~parseInt(s~\&\&~s.proactiveCustomSeconds,~10)\\
~~~~~~~~~~~~return~Number.isNaN(n)~?~0~:~Math.max(0,~Math.min(3600,~n))\\
~~~~~~~~\})(),\\
~~~~~~~~//~聊天气泡背景不透明度：钳制~0.2\textasciitilde{}1，1=完全不透明\\
~~~~~~~~bubbleOpacity:~(()~=>~\{\\
~~~~~~~~~~~~const~n~=~parseFloat(s~\&\&~s.bubbleOpacity)\\
~~~~~~~~~~~~return~Number.isFinite(n)~?~Math.max(0.2,~Math.min(1,~n))~:~1\\
~~~~~~~~\})()\\
~~~~\}\\
\}\\
\mbox{}\\
/**~保存设置到当前会话（每个对话独享一份设置，设置面板与之同步）~*/\\
export~function~saveSettings(s)~\{\\
~~~~setConversationSettingsRaw(normalizeSettings(s))\\
\}\\
\mbox{}\\
/**\\
~*~当前会话生效设置~=~会话设置快照（无快照回退全局；旧数据兼容仅人格子集的快照）。\\
~*~每个会话独立一套完整设置，切换会话即切换设置。\\
~*/\\
export~function~getConversationSettings()~\{\\
~~~~const~s~=~getSettings()\\
~~~~const~raw~=~getConversationSettingsRaw()\\
~~~~if~(raw)~\{\\
~~~~~~~~const~merged~=~\{~...s,~...raw~\}\\
~~~~~~~~merged.temperature~=~Number.isFinite(parseFloat(raw.temperature))~?~parseFloat(raw.temperatu\\
re)~:~s.temperature\\
~~~~~~~~merged.compressInterval~=~parseInt(raw.compressInterval,~10)~||~0\\
~~~~~~~~merged.maxRequestAttempts~=~Math.max(1,~Math.min(20,~parseInt(raw.maxRequestAttempts,~10)~||
}

\newpage

{\ttfamily\footnotesize\raggedright
~5))\\
~~~~~~~~merged.ttsEnabled~=~raw.ttsEnabled~===~undefined~?~s.ttsEnabled~:~raw.ttsEnabled~!==~false\\
~~~~~~~~merged.personaName~=~getPersonaName(merged.personalityId~||~s.personalityId)\\
~~~~~~~~return~merged\\
~~~~\}\\
~~~~//~旧数据：会话仅存人格子集（personalityId/customPrompt/timeMode）\\
~~~~const~p~=~getConversationPersonality()\\
~~~~return~\{\\
~~~~~~~~...s,\\
~~~~~~~~personalityId:~(p~\&\&~p.personalityId)~||~s.personalityId,\\
~~~~~~~~customPrompt:~p~\&\&~p.customPrompt~!==~undefined~?~p.customPrompt~:~s.customPrompt,\\
~~~~~~~~timeMode:~(p~\&\&~p.timeMode)~||~s.timeMode,\\
~~~~~~~~personaName:~getPersonaName((p~\&\&~p.personalityId)~||~s.personalityId)\\
~~~~\}\\
\}\\
\mbox{}\\
/**\\
~*~保存当前会话的人格设置（写入会话设置快照，仅作用于当前对话，不影响其他会话）。\\
~*~timeMode~可选：不传则沿用当前会话生效的模式，保持会话内的情景时间设置不漂移。\\
~*/\\
export~function~saveConversationPersonality(personalityId,~customPrompt,~timeMode)~\{\\
~~~~const~s~=~getConversationSettings()\\
~~~~setConversationSettingsRaw(\\
~~~~~~~~normalizeSettings(\{\\
~~~~~~~~~~~~...s,\\
~~~~~~~~~~~~personalityId,\\
~~~~~~~~~~~~customPrompt:~personalityId~===~'custom'~?~(customPrompt~||~'')~:~'',\\
~~~~~~~~~~~~timeMode:~timeMode~!==~undefined~?~timeMode~:~s.timeMode\\
~~~~~~~~\})\\
~~~~)\\
~~~~addLog('info',~'会话人格',~`\$\{getPersonaName(personalityId)\}\$\{personalityId~===~'custom'~?~'（自\\
定义）'~:~''\}`)\\
\}\\
\mbox{}\\
\mbox{}\\
//~================~FILE:~utils/chat-state.js~================\\
/**\\
~*~聊天服务共享状态~----~供~chat.js~/~chat-conversations.js~等模块共享的可变状态：\\
~*~-~memoryStore：记忆单例（会话切换重置~/~消息发送检索与入库~/~维护）\\
~*~-~lastRequest：最近一次请求内容缓存（"重新生成"直接重发，不重复落库用户消息）\\
~*~独立成模块以消除~chat.js~与~chat-conversations.js~之间的循环依赖。\\
~*/\\
\mbox{}\\
import~\{~MemoryStore~\}~from~'./memory.js'\\
\mbox{}\\
/**~记忆单例~*/\\
export~const~memoryStore~=~new~MemoryStore()\\
\mbox{}\\
/**~最近一次请求内容缓存（跨模块读写）~*/\\
export~const~lastRequest~=~\{~value:~''~\}
}

\newpage

{\ttfamily\footnotesize\raggedright
\mbox{}\\
\mbox{}\\
//~================~FILE:~utils/chat.js~================\\
/**\\
~*~聊天服务~----~门面~+~消息域：\\
~*~对外统一导出（其余域拆分至~chat-settings~/~chat-conversations~/~chat-compress），\\
~*~本文件承载发送主链路：检索记忆~->~组装~system~->~调用~LLM~->~落库~->~维护，\\
~*~以及"重新生成"的撤回逻辑与回复格式校验。\\
~*/\\
\mbox{}\\
import~\{\\
~~~~initStorage,\\
~~~~getChatRows,\\
~~~~addChatRow,\\
~~~~truncateChat,\\
~~~~getSetting,\\
~~~~setSetting,\\
~~~~getSceneHistory,\\
~~~~setScene,\\
~~~~truncateSceneHistory,\\
~~~~getConversationCompression\\
\}~from~'./storage.js'\\
import~\{~memoryStore,~lastRequest~\}~from~'./chat-state.js'\\
import~\{\\
~~~~personaName,\\
~~~~getSettings,\\
~~~~saveSettings,\\
~~~~getConversationSettings,\\
~~~~saveConversationPersonality\\
\}~from~'./chat-settings.js'\\
import~\{\\
~~~~getHistoryForUI,\\
~~~~clearConversation,\\
~~~~listConversations,\\
~~~~activeConversationId,\\
~~~~startNewConversation,\\
~~~~openConversation,\\
~~~~removeConversation,\\
~~~~copyConversationToNew,\\
~~~~copyMemoriesToNew\\
\}~from~'./chat-conversations.js'\\
import~\{~compressContext,~maybeCompress~\}~from~'./chat-compress.js'\\
import~\{~parseMemoryLine~\}~from~'./memory.js'\\
import~\{~buildSystemPrompt,~buildNowText,~getPersonalityById~\}~from~'./prompts.js'\\
import~\{~chatCompletion~\}~from~'./llm.js'\\
import~\{~addLog~\}~from~'./log.js'\\
import~\{~getEmojiMap,~emojiListForPrompt,~extractEmojiNames~\}~from~'./emojis.js'\\
\mbox{}\\
//~----------~对外统一导出（门面）~----------\\
//~chat.js~保持为唯一入口，各调用方（页面/工具/测试）的导入路径不变
}

\newpage

{\ttfamily\footnotesize\raggedright
\mbox{}\\
export~\{\\
~~~~memoryStore,\\
~~~~personaName,\\
~~~~getSettings,\\
~~~~saveSettings,\\
~~~~getConversationSettings,\\
~~~~saveConversationPersonality,\\
~~~~getHistoryForUI,\\
~~~~clearConversation,\\
~~~~listConversations,\\
~~~~activeConversationId,\\
~~~~startNewConversation,\\
~~~~openConversation,\\
~~~~removeConversation,\\
~~~~copyConversationToNew,\\
~~~~copyMemoriesToNew,\\
~~~~compressContext,\\
~~~~maybeCompress\\
\}\\
\mbox{}\\
//~----------~消息域~----------\\
\mbox{}\\
const~HISTORY\_ENTRIES~=~15~//~每轮注入的对话历史条数（上限）\\
const~MAINTENANCE\_INTERVAL\_MS~=~30~*~60~*~1000~//~维护节流间隔（30~分钟）\\
\mbox{}\\
/**~初始化存储（幂等，App.vue~onLaunch~调用）~*/\\
export~function~initChatService()~\{\\
~~~~initStorage()\\
~~~~memoryStore.maintenance()\\
~~~~setSetting('lastMaintenance',~new~Date().toISOString())\\
\}\\
\mbox{}\\
/**\\
~*~节流维护：距离上次维护超过间隔才执行（App~进入前台~/~页面~onShow~调用）。\\
~*~对话后的主动维护仍走~memoryStore.maintenance()，不受此限制。\\
~*/\\
export~function~maybeMaintenance()~\{\\
~~~~const~last~=~getSetting('lastMaintenance',~'')\\
~~~~if~(last)~\{\\
~~~~~~~~const~t~=~Date.parse(last)\\
~~~~~~~~if~(!Number.isNaN(t)~\&\&~Date.now()~-~t~<~MAINTENANCE\_INTERVAL\_MS)~return\\
~~~~\}\\
~~~~memoryStore.maintenance()\\
~~~~setSetting('lastMaintenance',~new~Date().toISOString())\\
\}\\
\mbox{}\\
/**\\
~*~用户点"重新生成"时调用：移除待重发请求对应的最后一条助手回复（撤回其情景与记忆）。\\
~*~返回待重发内容----优先取最近一次请求缓存，无缓存（如重启后）回退到当前会话最后一条用户消息；
}

\newpage

{\ttfamily\footnotesize\raggedright
~*~若无可重发内容则返回空串。用户消息保留在存储中，重发（persistUser:false）不再重复落库。\\
~*/\\
export~function~popLastAssistant()~\{\\
~~~~const~rows~=~getChatRows()\\
~~~~//~待重发内容：最近一次请求缓存优先；无缓存（如重启后）回退到最后一条用户消息\\
~~~~let~req~=~lastRequest.value\\
~~~~if~(!req)~\{\\
~~~~~~~~for~(let~i~=~rows.length~-~1;~i~>=~0;~i--)~\{\\
~~~~~~~~~~~~if~(rows[i].role~===~'user')~\{\\
~~~~~~~~~~~~~~~~req~=~rows[i].content\\
~~~~~~~~~~~~~~~~break\\
~~~~~~~~~~~~\}\\
~~~~~~~~\}\\
~~~~\}\\
~~~~if~(!req)~return~''\\
~~~~//~定位该请求对应的最后一条用户消息（从后往前匹配，取最近一次），\\
~~~~//~若其后有~assistant~回复则移除并撤回其情景/记忆（拟真模式~burst~为连续多条~assistant~行，\\
~~~~//~一并移除，rollback~挂在最后一行）；发送失败时无回复可移，仅重发\\
~~~~for~(let~i~=~rows.length~-~1;~i~>=~0;~i--)~\{\\
~~~~~~~~if~(rows[i].role~===~'user'~\&\&~rows[i].content~===~req)~\{\\
~~~~~~~~~~~~let~last~=~i\\
~~~~~~~~~~~~while~(last~+~1~<~rows.length~\&\&~rows[last~+~1].role~===~'assistant')~last++\\
~~~~~~~~~~~~if~(last~>~i)~\{\\
~~~~~~~~~~~~~~~~rollbackAssistantEffects(rows[last])\\
~~~~~~~~~~~~~~~~truncateChat(i~+~1)\\
~~~~~~~~~~~~\}\\
~~~~~~~~~~~~break\\
~~~~~~~~\}\\
~~~~\}\\
~~~~addLog('info',~'重新生成',~`重发最近一次请求：\$\{req.slice(0,~20)\}`)\\
~~~~return~req\\
\}\\
\mbox{}\\
/**~撤回一条响应记录附带的情景与记忆变更（重新生成前调用）~*/\\
function~rollbackAssistantEffects(row)~\{\\
~~~~const~rb~=~row~\&\&~row.rollback\\
~~~~if~(!rb)~return\\
~~~~//~情景：截断到该响应之前的情景历史长度，撤掉本响应新增的情景\\
~~~~if~(typeof~rb.sceneLenBefore~===~'number')~\{\\
~~~~~~~~truncateSceneHistory(rb.sceneLenBefore)\\
~~~~\}\\
~~~~//~记忆：按写入顺序逆序撤销（新增->删除~/~修改->恢复原值~/~删除->重新插入）\\
~~~~const~undos~=~rb.memoryUndos\\
~~~~if~(Array.isArray(undos)~\&\&~undos.length)~\{\\
~~~~~~~~for~(let~i~=~undos.length~-~1;~i~>=~0;~i--)~memoryStore.applyUndo(undos[i])\\
~~~~\}\\
~~~~addLog('info',~'撤回响应情景与记忆',~`情景截断到~\$\{rb.sceneLenBefore\}，撤销记忆~\$\{(undos~\&\&~undo\\
s.length)~||~0\}~条`)\\
\}\\
\mbox{}
}

\newpage

{\ttfamily\footnotesize\raggedright
/**\\
~*~发送一条用户消息\\
~*~@param~\{string\}~userText\\
~*~@param~\{\{persistUser?:~boolean\}\}~[opts]~persistUser=false~表示"重新生成"重发：\\
~*~~~~~~~~用户消息已在首次发送时落库，本次只更新回复，避免重复记录同一句话\\
~*~@returns~\{Promise<\{reply:string,~saved:number\}>\}~reply~为清理掉~Memory~行后的回复\\
~*/\\
export~async~function~sendMessage(userText,~opts~=~\{\})~\{\\
~~~~//~设置全部取当前会话（每个对话独享一份设置：API~配置~/~人格~/~时间模式等）\\
~~~~const~persistUser~=~opts.persistUser~!==~false\\
~~~~const~s~=~getConversationSettings()\\
~~~~if~(!s.apiKey)~throw~new~Error('请先在「设置」中填写~API~Key')\\
~~~~if~(!s.baseUrl)~throw~new~Error('请先在「设置」中填写接口地址')\\
~~~~if~(!s.model)~throw~new~Error('请先在「设置」中填写模型名称')\\
~~~~//~拟真聊天：该会话开启后，所有回复均受"每条~<=1~句、连续表情~<=2"约束，并按换行拆分为多条气泡\\
~~~~const~proactive~=~!!s.proactiveEnabled\\
\mbox{}\\
~~~~//~缓存最近一次请求：聊天页"重新生成"直接重发该内容（重发时~persistUser:false，不重复落库用户消\\
息）\\
~~~~lastRequest.value~=~userText\\
\mbox{}\\
~~~~//~记录发送前的情景历史长度：若本响应更新了情景，重新生成时据此撤回\\
~~~~const~sceneLenBefore~=~getSceneHistory().length\\
\mbox{}\\
~~~~//~1.~检索记忆（全量~L1~+~新鲜槽~+~MMR~多样性槽）\\
~~~~const~memoryText~=~memoryStore.retrieveContext(userText)\\
~~~~addLog('info',~'发送消息',~userText.slice(0,~40))\\
\mbox{}\\
~~~~//~2.~组装~system~prompt（人格~+~规则~+~记忆指南~+~情景指南~+~当前状态~+~记忆~+~L1~容量~+~表情包\\
清单）\\
~~~~const~personalityPrompt~=\\
~~~~~~~~s.personalityId~===~'custom'~?~s.customPrompt~:~getPersonalityById(s.personalityId).prompt\\
~~~~const~system~=~buildSystemPrompt(\\
~~~~~~~~personalityPrompt~||~'你是友好的聊天伙伴。',\\
~~~~~~~~memoryText,\\
~~~~~~~~getSceneHistory(),\\
~~~~~~~~s.timeMode~===~'virtual'~?~''~:~buildNowText(),\\
~~~~~~~~//~表情包开关：禁用时不传清单（buildSystemPrompt~不注入表情包规则，LLM~不会主动使用表情）\\
~~~~~~~~s.emojiEnabled~?~emojiListForPrompt()~:~[],\\
~~~~~~~~//~L1~容量状态：引导~LLM~在~L1~满员时先逐字删除/修改旧~L1~再新增\\
~~~~~~~~memoryStore.l1Usage(),\\
~~~~~~~~//~拟真聊天：注入每条消息~<=1~句的规则（开启时生效）\\
~~~~~~~~proactive\\
~~~~)\\
\mbox{}\\
~~~~//~3-5.~请求~->~格式校验~->~解析：回复格式不合格时自动重试（上限~maxRequestAttempts）。\\
~~~~//~每次尝试都基于同一份历史重新组装；校验通过才写入~Scene/Memory~并落库，失败尝试无副作用。\\
~~~~const~maxAttempts~=~Math.max(1,~Math.min(20,~parseInt(s.maxRequestAttempts,~10)~||~5))\\
~~~~let~result~=~null\\
~~~~let~lastReason~=~''
}

\newpage

{\ttfamily\footnotesize\raggedright
~~~~//~任一尝试失败（网络错误~/~格式校验达上限）都只记录用户请求、不落库错误回复，\\
~~~~//~聊天页"重新生成"基于最近一次请求缓存直接重发\\
~~~~try~\{\\
~~~~~~~~for~(let~attempt~=~1;~attempt~<=~maxAttempts;~attempt++)~\{\\
~~~~~~~~~~~~//~组装~messages：system~+~未压缩历史~+~当前用户消息\\
~~~~~~~~~~~~//~压缩概要并入首条~system~末尾，避免出现第二条~system~消息----Ollama~等模板要求\\
~~~~~~~~~~~~//~system~必须在~messages~最前且只能一条，多条~system~会抛~"System~message~must~be~at~th\\
e~beginning"\\
~~~~~~~~~~~~const~\{~summary,~compressedUntil~\}~=~getConversationCompression()\\
~~~~~~~~~~~~const~messages~=~[\\
~~~~~~~~~~~~~~~~\{~role:~'system',~content:~summary~?~`\$\{system\}\textbackslash{}n\textbackslash{}n此前对话概要：\$\{summary\}`~:~syste\\
m~\}\\
~~~~~~~~~~~~]\\
~~~~~~~~~~~~for~(const~h~of~getChatRows().slice(compressedUntil).slice(-HISTORY\_ENTRIES))~\{\\
~~~~~~~~~~~~~~~~messages.push(\{~role:~h.role,~content:~h.content~\})\\
~~~~~~~~~~~~\}\\
~~~~~~~~~~~~messages.push(\{~role:~'user',~content:~userText~\})\\
\mbox{}\\
~~~~~~~~~~~~const~reply~=~await~chatCompletion(\{\\
~~~~~~~~~~~~~~~~baseUrl:~s.baseUrl,\\
~~~~~~~~~~~~~~~~apiKey:~s.apiKey,\\
~~~~~~~~~~~~~~~~model:~s.model,\\
~~~~~~~~~~~~~~~~messages,\\
~~~~~~~~~~~~~~~~temperature:~s.temperature,\\
~~~~~~~~~~~~~~~~reasoningEffort:~s.reasoningEffort~//~默认关闭思考（none），本地思考模型（如~Qwen3.5\\
）默认思考会占满输出~token~导致回复为空\\
~~~~~~~~~~~~\})\\
\mbox{}\\
~~~~~~~~~~~~const~parsed~=~parseAndValidateReply(reply.text,~\{~proactive~\})\\
~~~~~~~~~~~~if~(parsed.ok)~\{\\
~~~~~~~~~~~~~~~~result~=~parsed\\
~~~~~~~~~~~~~~~~break\\
~~~~~~~~~~~~\}\\
~~~~~~~~~~~~lastReason~=~parsed.reason\\
~~~~~~~~~~~~addLog('info',~'回复格式不合格，自动重试',~`第~\$\{attempt\}/\$\{maxAttempts\}~次：\$\{lastReaso\\
n\}`)\\
~~~~~~~~\}\\
~~~~~~~~if~(!result)~\{\\
~~~~~~~~~~~~addLog('err',~'回复格式不合格（已达请求上限）',~`最大请求次数~\$\{maxAttempts\}：\$\{lastReas\\
on\}`)\\
~~~~~~~~~~~~throw~new~Error(`模型回复格式连续~\$\{maxAttempts\}~次不合格（\$\{lastReason\}）`)\\
~~~~~~~~\}\\
~~~~\}~catch~(e)~\{\\
~~~~~~~~//~发送失败：只落库用户请求（不落库错误回复），供"重新生成"定位本次请求\\
~~~~~~~~if~(persistUser)~addChatRow('user',~userText)\\
~~~~~~~~throw~e\\
~~~~\}\\
\mbox{}\\
~~~~//~6.~落库对话：落库清理后的文本，Scene/Memory~标记不再混入历史----否则从其它页切回时\\
~~~~//~重新加载历史会把标记显示在气泡里（此前~bug），也会污染上下文压缩。
}

\newpage

{\ttfamily\footnotesize\raggedright
~~~~//~assistant~行附带~rollback~信息（响应前的情景长度~+~本响应的记忆撤销清单），\\
~~~~//~重新生成时据此撤回该响应记录的情景与记忆。\\
~~~~//~拟真模式：按换行拆分多条落库（每条~<=1~句一条气泡），rollback~挂在最后一行\\
~~~~if~(persistUser)~addChatRow('user',~userText)\\
~~~~const~cleanReply~=~result.cleanReply~||~'...'\\
~~~~const~rollback~=~\{~sceneLenBefore,~memoryUndos:~result.memoryUndos~||~[]~\}\\
~~~~if~(proactive)~\{\\
~~~~~~~~const~lines~=~cleanReply.split('\textbackslash{}n').map((l)~=>~l.trim()).filter(Boolean)\\
~~~~~~~~if~(!lines.length)~lines.push('...')\\
~~~~~~~~lines.forEach((line,~idx)~=>~\{\\
~~~~~~~~~~~~addChatRow('assistant',~line,~idx~===~lines.length~-~1~?~\{~rollback~\}~:~undefined)\\
~~~~~~~~\})\\
~~~~\}~else~\{\\
~~~~~~~~addChatRow('assistant',~cleanReply,~\{~rollback~\})\\
~~~~\}\\
~~~~if~(result.newScene)~\{\\
~~~~~~~~setScene(result.newScene)\\
~~~~~~~~addLog('info',~'情景更新',~result.newScene)\\
~~~~\}\\
~~~~if~(result.saved~>~0)~addLog('info',~`记忆入库~\$\{result.saved\}~条`)\\
\mbox{}\\
~~~~//~7.~定期维护（L3~过期清理~/~降级~/~容量控制）\\
~~~~memoryStore.maintenance()\\
\mbox{}\\
~~~~//~8.~自动压缩检查（异步执行，不阻塞回复展示）\\
~~~~maybeCompress().catch(()~=>~\{~\})\\
\mbox{}\\
~~~~return~\{~reply:~cleanReply,~saved:~result.saved,~burst:~proactive~?~cleanReply.split('\textbackslash{}n').map((\\
l)~=>~l.trim()).filter(Boolean)~:~undefined~\}\\
\}\\
\mbox{}\\
/**\\
~*~统计文本句子数：按句末标点（。！？!?...）计，连续标点（如"哈哈。。"）合并为一处，\\
~*~用于拟真聊天"每条消息~<=1~句"的校验。\\
~*~@param~\{string\}~text\\
~*~@returns~\{number\}\\
~*/\\
export~function~countSentences(text)~\{\\
~~~~const~t~=~String(text~||~'').trim()\\
~~~~if~(!t)~return~0\\
~~~~let~runs~=~0\\
~~~~let~inRun~=~false\\
~~~~for~(const~ch~of~t)~\{\\
~~~~~~~~if~('。！？!?...'.includes(ch))~\{\\
~~~~~~~~~~~~if~(!inRun)~\{\\
~~~~~~~~~~~~~~~~runs++\\
~~~~~~~~~~~~~~~~inRun~=~true\\
~~~~~~~~~~~~\}\\
~~~~~~~~\}~else~\{\\
~~~~~~~~~~~~inRun~=~false
}

\newpage

{\ttfamily\footnotesize\raggedright
~~~~~~~~\}\\
~~~~\}\\
~~~~return~runs\\
\}\\
\mbox{}\\
/**\\
~*~校验~LLM~回复格式并解析~Scene/Memory（校验通过才写库，供~sendMessage~重试使用）：\\
~*~-~回复必须包含对话文本（仅~Scene/Memory~标记视为不合格）\\
~*~-~Scene~行必须有内容；Memory~行必须能被~parseMemoryLine~解析\\
~*~-~拟真模式（proactive:true）：每条文本行~<=1~句、连续表情~<=2，不合格触发重试\\
~*~@param~\{string\}~text~LLM~原始回复\\
~*~@param~\{\{proactive?:~boolean\}\}~[opts]\\
~*~@returns~\{\{ok:boolean,~reason?:string,~saved?:number,~newScene?:string|null,~cleanReply?:string,~\\
memoryUndos?:Array\}\}\\
~*/\\
export~function~parseAndValidateReply(text,~opts)~\{\\
~~~~const~proactive~=~!!(opts~\&\&~opts.proactive)\\
~~~~//~表情名校验：回复中的~\$名\$~必须存在于表情清单，否则判定格式不合格触发重试\\
~~~~//~（避免~LLM~编造清单外的名字导致聊天页渲染出"坏图"）\\
~~~~const~emojiMap~=~getEmojiMap()\\
~~~~for~(const~name~of~extractEmojiNames(text))~\{\\
~~~~~~~~if~(!emojiMap[name])~return~\{~ok:~false,~reason:~`回复包含未知表情名：\$\{name\}`~\}\\
~~~~\}\\
~~~~const~lines~=~String(text~||~'').split('\textbackslash{}n')\\
~~~~//~第一步：格式校验（不写入）\\
~~~~let~newScene~=~null\\
~~~~let~hasContent~=~false\\
~~~~const~cleanLines~=~[]\\
~~~~for~(const~line~of~lines)~\{\\
~~~~~~~~if~(/\textasciicircum{}\textbackslash{}s*Scene\textbackslash{}s*[:：]/i.test(line))~\{\\
~~~~~~~~~~~~newScene~=~line\\
~~~~~~~~~~~~~~~~.replace(/\textasciicircum{}\textbackslash{}s*Scene\textbackslash{}s*[:：]\textbackslash{}s*/i,~'')\\
~~~~~~~~~~~~~~~~.replace(/\textasciicircum{}["'""'']+|["'""'']+\$/g,~'')~//~去掉~LLM~可能加上的引号\\
~~~~~~~~~~~~~~~~.replace(/[。！!]\$/,~'')~~~~~~~~~~~~~~//~去掉末尾标点\\
~~~~~~~~~~~~~~~~.trim()\\
~~~~~~~~~~~~if~(!newScene)~return~\{~ok:~false,~reason:~'Scene~行内容为空'~\}\\
~~~~~~~~~~~~continue\\
~~~~~~~~\}\\
~~~~~~~~if~(/\textasciicircum{}\textbackslash{}s*Memory\textbackslash{}s*[:：]/i.test(line))~\{\\
~~~~~~~~~~~~if~(!parseMemoryLine(line))~return~\{~ok:~false,~reason:~`Memory~行格式非法：\$\{line.trim(\\
).slice(0,~24)\}`~\}\\
~~~~~~~~~~~~continue\\
~~~~~~~~\}\\
~~~~~~~~//~行内~Scene/Memory~标记（如~"正文~[Scene:~xxx]"）：Scene/Memory~必须独立成行，\\
~~~~~~~~//~行首形式已在上方分支处理，走到这里仍含~Scene:/Memory:~即未独立成行~->~判定格式非法\\
~~~~~~~~if~(/Scene\textbackslash{}s*[:：]/i.test(line)~||~/Memory\textbackslash{}s*[:：]/i.test(line))~\{\\
~~~~~~~~~~~~return~\{~ok:~false,~reason:~`Scene/Memory~标记未独立成行：\$\{line.trim().slice(0,~28)\}`~\}\\
~~~~~~~~\}\\
~~~~~~~~//~疑似~Memory~行：含~|~keywords:/|~importance:/|~level:~分段结构却缺少~Memory:~前缀\\
~~~~~~~~//~（如~"[OK]~更新：3.~xxx~|~keywords:..~|~importance:..~|~level:.."），判定格式非法并触发重
}

\newpage

{\ttfamily\footnotesize\raggedright
试\\
~~~~~~~~if~(/\textbackslash{}|\textbackslash{}s*(?:keywords|importance|level):/i.test(line))~\{\\
~~~~~~~~~~~~return~\{~ok:~false,~reason:~`疑似~Memory~行缺少~Memory:~前缀：\$\{line.trim().slice(0,~28)\\
\}`~\}\\
~~~~~~~~\}\\
~~~~~~~~//~拟真模式：每条消息（一行）<=1~句、连续表情~<=2，超出即判定格式不合格触发重试\\
~~~~~~~~if~(proactive)~\{\\
~~~~~~~~~~~~if~(countSentences(line)~>~1)~\{\\
~~~~~~~~~~~~~~~~return~\{~ok:~false,~reason:~`拟真模式消息超过一句话：\$\{line.trim().slice(0,~20)\}`~\}\\
~~~~~~~~~~~~\}\\
~~~~~~~~~~~~if~(extractEmojiNames(line).length~>~2)~\{\\
~~~~~~~~~~~~~~~~return~\{~ok:~false,~reason:~`拟真模式连续表情超过~2~个：\$\{line.trim().slice(0,~20)\}`\\
~\}\\
~~~~~~~~~~~~\}\\
~~~~~~~~\}\\
~~~~~~~~cleanLines.push(line)\\
~~~~~~~~if~(line.trim())~hasContent~=~true\\
~~~~\}\\
~~~~if~(!hasContent)~return~\{~ok:~false,~reason:~'回复仅含~Scene/Memory~标记，无对话文本'~\}\\
~~~~const~cleanReply~=~cleanLines.join('\textbackslash{}n').trim()\\
\mbox{}\\
~~~~//~第二步：校验通过，执行写入（Scene~更新情景~/~Memory~入库），并收集撤销信息供重新生成回滚\\
~~~~let~saved~=~0\\
~~~~const~memoryUndos~=~[]\\
~~~~for~(const~line~of~lines)~\{\\
~~~~~~~~if~(/\textasciicircum{}\textbackslash{}s*Memory\textbackslash{}s*[:：]/i.test(line))~\{\\
~~~~~~~~~~~~const~undo~=~memoryStore.saveFromLine(line)\\
~~~~~~~~~~~~if~(undo)~\{\\
~~~~~~~~~~~~~~~~saved++\\
~~~~~~~~~~~~~~~~memoryUndos.push(undo)\\
~~~~~~~~~~~~\}\\
~~~~~~~~\}\\
~~~~\}\\
~~~~return~\{~ok:~true,~saved,~newScene,~cleanReply,~memoryUndos~\}\\
\}\\
\mbox{}\\
\mbox{}\\
//~================~FILE:~utils/emojis.js~================\\
/**\\
~*~表情包数据层~----~全局共享的表情列表（不随会话切换）\\
~*\\
~*~每个表情~\{~id,~name,~src,~created\_at~\}：\\
~*~-~name~表情名：上传时必填、全局唯一、<=20~字、不含~\$/换行；消息中用~\$表情名\$~引用\\
~*~-~src~图片：App~/~微信小程序为~uni.saveFile~持久化路径，H5~为压缩后的~PNG~base64\\
~*\\
~*~存储键~chabot\_emojis（uni~storage，跨端可用）。表情数量与图片大小设上限，\\
~*~避免~H5~localStorage~配额（约~5MB）被大图打满。\\
~*/\\
\mbox{}\\
const~KEY\_EMOJIS~=~'chabot\_emojis'
}

\newpage

{\ttfamily\footnotesize\raggedright
\mbox{}\\
export~const~EMOJI\_MAX\_COUNT~=~50~//~表情数量上限\\
export~const~EMOJI\_MAX\_NAME\_LEN~=~20~//~表情名长度上限（字），与~\$名\$~占位解析的正则上限一致\\
const~EMOJI\_IMG\_MAX\_WIDTH~=~300~//~H5~压缩后的图片最大宽度（px）\\
const~EMOJI\_IMG\_MAX\_BYTES~=~300~*~1024~//~H5~base64~单图大小上限（约~300KB）\\
\mbox{}\\
let~\_emojis~=~[]~//~内存缓存（活数据）\\
let~\_loaded~=~false\\
\mbox{}\\
/**~首次访问时从存储加载（幂等）~*/\\
function~\_ensureLoaded()~\{\\
~~~~if~(\_loaded)~return\\
~~~~\_loaded~=~true\\
~~~~if~(typeof~uni~!==~'undefined'~\&\&~uni.getStorageSync)~\{\\
~~~~~~~~const~v~=~uni.getStorageSync(KEY\_EMOJIS)\\
~~~~~~~~\_emojis~=~Array.isArray(v)~?~v.filter((e)~=>~e~\&\&~typeof~e.name~===~'string'~\&\&~e.name~\&\&~e.\\
src)~:~[]\\
~~~~\}~else~\{\\
~~~~~~~~\_emojis~=~[]\\
~~~~\}\\
\}\\
\mbox{}\\
function~\_persist()~\{\\
~~~~if~(typeof~uni~!==~'undefined'~\&\&~uni.setStorageSync)~\{\\
~~~~~~~~try~\{\\
~~~~~~~~~~~~uni.setStorageSync(KEY\_EMOJIS,~\_emojis)\\
~~~~~~~~\}~catch~(e)~\{\\
~~~~~~~~~~~~console.error('[emojis]~保存失败:',~e)\\
~~~~~~~~\}\\
~~~~\}\\
\}\\
\mbox{}\\
/**~表情列表（副本）~*/\\
export~function~getEmojis()~\{\\
~~~~\_ensureLoaded()\\
~~~~return~\_emojis.map((e)~=>~(\{~...e~\}))\\
\}\\
\mbox{}\\
/**~表情名~->~图片~映射（渲染与校验用）~*/\\
export~function~getEmojiMap()~\{\\
~~~~\_ensureLoaded()\\
~~~~const~map~=~\{\}\\
~~~~for~(const~e~of~\_emojis)~map[e.name]~=~e.src\\
~~~~return~map\\
\}\\
\mbox{}\\
/**~表情名清单（供~LLM~提示词注入）~*/\\
export~function~emojiListForPrompt()~\{\\
~~~~\_ensureLoaded()\\
~~~~return~\_emojis.map((e)~=>~e.name)
}

\newpage

{\ttfamily\footnotesize\raggedright
\}\\
\mbox{}\\
/**\\
~*~校验表情名：必填~/~<=20~字~/~不含~\$~与换行~/~全局唯一（改名时排除自身）\\
~*~@param~\{string\}~name\\
~*~@param~\{string\}~[excludeId]~排除的表情~id（改名时传入自身）\\
~*~@returns~\{\{ok:boolean,~reason?:string\}\}\\
~*/\\
export~function~validateEmojiName(name,~excludeId)~\{\\
~~~~const~n~=~String(name~||~'').trim()\\
~~~~if~(!n)~return~\{~ok:~false,~reason:~'请填写表情名'~\}\\
~~~~if~(n.length~>~EMOJI\_MAX\_NAME\_LEN)~return~\{~ok:~false,~reason:~`表情名不能超过~\$\{EMOJI\_MAX\_NAME\_\\
LEN\}~字`~\}\\
~~~~if~(n.indexOf('\$')~!==~-1~||~/[\textbackslash{}r\textbackslash{}n]/.test(n))~return~\{~ok:~false,~reason:~'表情名不能包含~\$~或\\
换行'~\}\\
~~~~\_ensureLoaded()\\
~~~~for~(const~e~of~\_emojis)~\{\\
~~~~~~~~if~(e.name~===~n~\&\&~e.id~!==~excludeId)~return~\{~ok:~false,~reason:~`表情名「\$\{n\}」已存在`~\}\\
~~~~\}\\
~~~~return~\{~ok:~true~\}\\
\}\\
\mbox{}\\
/**\\
~*~以已持久化的~src~直接新增表情（图片保存由~addEmoji~负责；测试/导入场景可直用）\\
~*~@param~\{string\}~name~表情名\\
~*~@param~\{string\}~src~图片路径~/~base64\\
~*~@returns~\{\{id:string,~name:string,~src:string,~created\_at:string\}\}~新增的表情副本\\
~*/\\
export~function~addEmojiData(name,~src)~\{\\
~~~~const~v~=~validateEmojiName(name)\\
~~~~if~(!v.ok)~throw~new~Error(v.reason)\\
~~~~\_ensureLoaded()\\
~~~~if~(\_emojis.length~>=~EMOJI\_MAX\_COUNT)~throw~new~Error(`表情数量已达上限（\$\{EMOJI\_MAX\_COUNT\}~个\\
）`)\\
~~~~if~(!src)~throw~new~Error('缺少表情图片')\\
~~~~const~item~=~\{\\
~~~~~~~~id:~'e'~+~Date.now().toString(36)~+~'\_'~+~Math.floor(Math.random()~*~1e4),\\
~~~~~~~~name:~name.trim(),\\
~~~~~~~~src,\\
~~~~~~~~created\_at:~new~Date().toISOString()\\
~~~~\}\\
~~~~\_emojis.push(item)\\
~~~~\_persist()\\
~~~~return~\{~...item~\}\\
\}\\
\mbox{}\\
/**~保存表情图片为持久化~src（App/小程序压缩后存文件，H5~转小尺寸~PNG~base64），失败返回空串~*/\\
async~function~\_saveEmojiImage(tempPath)~\{\\
~~~~let~saved~=~''\\
~~~~//~\#ifdef~APP-PLUS~||~MP-WEIXIN
}

\newpage

{\ttfamily\footnotesize\raggedright
~~~~//~上传时先压缩（限宽~+~quality），避免原图过大：消息列表渲染大量大图会明显拖慢~App~启动与滚动\\
~~~~//~注意：compressImage~输出为~jpg，PNG~透明底会变白，表情多为此可接受；H5~端仍走~PNG~base64~保留\\
透明\\
~~~~if~(typeof~uni.compressImage~===~'function')~\{\\
~~~~~~~~const~compressed~=~await~new~Promise((resolve)~=>~\{\\
~~~~~~~~~~~~uni.compressImage(\{\\
~~~~~~~~~~~~~~~~src:~tempPath,\\
~~~~~~~~~~~~~~~~quality:~80,\\
~~~~~~~~~~~~~~~~compressedWidth:~EMOJI\_IMG\_MAX\_WIDTH,\\
~~~~~~~~~~~~~~~~success:~(res)~=>~resolve(res.tempFilePath~||~''),\\
~~~~~~~~~~~~~~~~fail:~()~=>~resolve('')\\
~~~~~~~~~~~~\})\\
~~~~~~~~\})\\
~~~~~~~~if~(compressed)~tempPath~=~compressed\\
~~~~\}\\
~~~~saved~=~await~new~Promise((resolve)~=>~\{\\
~~~~~~~~uni.saveFile(\{\\
~~~~~~~~~~~~tempFilePath:~tempPath,\\
~~~~~~~~~~~~success:~(res)~=>~resolve(res.savedFilePath~||~''),\\
~~~~~~~~~~~~fail:~()~=>~resolve('')\\
~~~~~~~~\})\\
~~~~\})\\
~~~~//~\#endif\\
~~~~//~\#ifdef~H5\\
~~~~saved~=~await~\_h5EmojiToBase64(tempPath)\\
~~~~//~\#endif\\
~~~~return~saved\\
\}\\
\mbox{}\\
//~\#ifdef~H5\\
/**~H5~端：canvas~把图片压缩为小尺寸~PNG~base64（保留透明，限制宽度与体积）~*/\\
function~\_h5EmojiToBase64(tempPath)~\{\\
~~~~return~new~Promise((resolve)~=>~\{\\
~~~~~~~~const~img~=~new~Image()\\
~~~~~~~~img.onload~=~()~=>~\{\\
~~~~~~~~~~~~try~\{\\
~~~~~~~~~~~~~~~~let~w~=~img.naturalWidth~||~img.width\\
~~~~~~~~~~~~~~~~let~h~=~img.naturalHeight~||~img.height\\
~~~~~~~~~~~~~~~~if~(w~>~EMOJI\_IMG\_MAX\_WIDTH)~\{\\
~~~~~~~~~~~~~~~~~~~~h~=~Math.round((h~*~EMOJI\_IMG\_MAX\_WIDTH)~/~w)\\
~~~~~~~~~~~~~~~~~~~~w~=~EMOJI\_IMG\_MAX\_WIDTH\\
~~~~~~~~~~~~~~~~\}\\
~~~~~~~~~~~~~~~~const~canvas~=~document.createElement('canvas')\\
~~~~~~~~~~~~~~~~canvas.width~=~w\\
~~~~~~~~~~~~~~~~canvas.height~=~h\\
~~~~~~~~~~~~~~~~canvas.getContext('2d').drawImage(img,~0,~0,~w,~h)\\
~~~~~~~~~~~~~~~~const~dataUrl~=~canvas.toDataURL('image/png')\\
~~~~~~~~~~~~~~~~//~base64~长度~x~3/4~\textasciitilde{}~字节数，超限视为失败（提示重新选择更小图片）\\
~~~~~~~~~~~~~~~~if~(dataUrl.length~*~0.75~>~EMOJI\_IMG\_MAX\_BYTES)~return~resolve('')\\
~~~~~~~~~~~~~~~~resolve(dataUrl)
}

\newpage

{\ttfamily\footnotesize\raggedright
~~~~~~~~~~~~\}~catch~(e)~\{\\
~~~~~~~~~~~~~~~~console.error('[emojis]~图片压缩失败',~e)\\
~~~~~~~~~~~~~~~~resolve('')\\
~~~~~~~~~~~~\}\\
~~~~~~~~\}\\
~~~~~~~~img.onerror~=~()~=>~resolve('')\\
~~~~~~~~img.src~=~tempPath\\
~~~~\})\\
\}\\
//~\#endif\\
\mbox{}\\
/**\\
~*~选择图片后新增表情（校验~->~保存图片~->~落库）\\
~*~@param~\{string\}~tempPath~uni.chooseImage~返回的临时路径\\
~*~@param~\{string\}~name~表情名\\
~*~@returns~\{Promise<object>\}~新增的表情副本\\
~*/\\
export~async~function~addEmoji(tempPath,~name)~\{\\
~~~~const~v~=~validateEmojiName(name)\\
~~~~if~(!v.ok)~throw~new~Error(v.reason)\\
~~~~if~(!tempPath)~throw~new~Error('未选择图片')\\
~~~~const~src~=~await~\_saveEmojiImage(tempPath)\\
~~~~if~(!src)~throw~new~Error('图片保存失败，请换一张图片重试')\\
~~~~return~addEmojiData(name,~src)\\
\}\\
\mbox{}\\
/**~修改表情名（唯一性校验排除自身）~*/\\
export~function~renameEmoji(id,~newName)~\{\\
~~~~const~v~=~validateEmojiName(newName,~id)\\
~~~~if~(!v.ok)~throw~new~Error(v.reason)\\
~~~~\_ensureLoaded()\\
~~~~const~e~=~\_emojis.find((x)~=>~x.id~===~id)\\
~~~~if~(!e)~throw~new~Error('表情不存在')\\
~~~~e.name~=~newName.trim()\\
~~~~\_persist()\\
~~~~return~\{~...e~\}\\
\}\\
\mbox{}\\
/**~删除表情（顺带清理已保存的图片文件），不存在返回~false~*/\\
export~function~deleteEmoji(id)~\{\\
~~~~\_ensureLoaded()\\
~~~~const~idx~=~\_emojis.findIndex((e)~=>~e.id~===~id)\\
~~~~if~(idx~<~0)~return~false\\
~~~~const~[removed]~=~\_emojis.splice(idx,~1)\\
~~~~\_persist()\\
~~~~if~(removed.src)~\{\\
~~~~~~~~//~\#ifdef~APP-PLUS~||~MP-WEIXIN\\
~~~~~~~~uni.removeSavedFile(\{~filePath:~removed.src,~fail:~()~=>~\{~\}~\})\\
~~~~~~~~//~\#endif\\
~~~~\}
}

\newpage

{\ttfamily\footnotesize\raggedright
~~~~return~true\\
\}\\
\mbox{}\\
/**\\
~*~按新顺序重排表情列表（拖拽排序用），列表外的~id~自动忽略。\\
~*~@param~\{string[]\}~orderedIds~新顺序的表情~id~数组\\
~*/\\
export~function~reorderEmojis(orderedIds)~\{\\
~~~~\_ensureLoaded()\\
~~~~const~byId~=~new~Map(\_emojis.map((e)~=>~[e.id,~e]))\\
~~~~const~next~=~(Array.isArray(orderedIds)~?~orderedIds~:~[])\\
~~~~~~~~.map((id)~=>~byId.get(id))\\
~~~~~~~~.filter(Boolean)\\
~~~~//~补上未出现在新顺序里的表情（防御：保证列表不丢项）\\
~~~~for~(const~e~of~\_emojis)~\{\\
~~~~~~~~if~(!next.includes(e))~next.push(e)\\
~~~~\}\\
~~~~\_emojis~=~next\\
~~~~\_persist()\\
\}\\
\mbox{}\\
/**~新建正则：匹配~\$表情名\$~占位（名~1\textasciitilde{}20~字，不含~\$~与换行）~*/\\
function~\_emojiTokenRe()~\{\\
~~~~return~/\textbackslash{}\$([\textasciicircum{}\$\textbackslash{}n]\{1,20\})\textbackslash{}\$/g\\
\}\\
\mbox{}\\
/**~提取文本中所有~\$表情名\$~占位名（含清单外的未知名），供~LLM~回复校验使用~*/\\
export~function~extractEmojiNames(text)~\{\\
~~~~const~names~=~[]\\
~~~~const~re~=~\_emojiTokenRe()\\
~~~~let~m\\
~~~~while~((m~=~re.exec(String(text~||~'')))~!==~null)~names.push(m[1])\\
~~~~return~names\\
\}\\
\mbox{}\\
/**\\
~*~把消息文本按~\$表情名\$~拆分为段数组（纯函数，供消息渲染与测试）。\\
~*~表情包前后的空格/换行等纯空白段会被跳过，文本段去除首尾空白----\\
~*~避免模型在表情之间插入空格/换行时渲染出空消息气泡。\\
~*~@param~\{string\}~content~原始消息（可含~\$名\$~占位）\\
~*~@param~\{Object\}~[map]~表情名~->~图片（getEmojiMap()~的结果）\\
~*~@returns~\{Array<\{type:'text',~text:string\}|\{type:'emoji',~name:string,~src:string\}>\}\\
~*~~~~~~~~~~清单外的未知~\$名\$~原样保留在文本段中\\
~*/\\
export~function~splitEmojiText(content,~map)~\{\\
~~~~const~text~=~String(content~||~'')\\
~~~~const~out~=~[]\\
~~~~const~re~=~\_emojiTokenRe()\\
~~~~let~last~=~0\\
~~~~let~m
}

\newpage

{\ttfamily\footnotesize\raggedright
~~~~while~((m~=~re.exec(text))~!==~null)~\{\\
~~~~~~~~const~name~=~m[1]\\
~~~~~~~~const~seg~=~text.slice(last,~m.index)\\
~~~~~~~~//~表情前的文本段：纯空白跳过，非空白去除首尾空白\\
~~~~~~~~const~trimmed~=~seg.trim()\\
~~~~~~~~if~(trimmed)~out.push(\{~type:~'text',~text:~trimmed~\})\\
~~~~~~~~const~src~=~map~\&\&~map[name]\\
~~~~~~~~out.push(src~?~\{~type:~'emoji',~name,~src~\}~:~\{~type:~'text',~text:~m[0]~\})\\
~~~~~~~~last~=~re.lastIndex\\
~~~~\}\\
~~~~if~(last~<~text.length)~\{\\
~~~~~~~~const~trimmed~=~text.slice(last).trim()\\
~~~~~~~~if~(trimmed)~out.push(\{~type:~'text',~text:~trimmed~\})\\
~~~~\}\\
~~~~return~out\\
\}\\
\mbox{}\\
\mbox{}\\
//~================~FILE:~utils/export.js~================\\
/**\\
~*~聊天记录导出~----~把当前会话聊天记录组装为纯文本，按平台导出一键~.txt：\\
~*~-~H5：Blob~直接触发浏览器下载（带~BOM~防中文乱码）\\
~*~-~App：plus.io~写入应用文档目录（\_doc），并尝试系统打开\\
~*~-~微信小程序：小程序无法直接导出~txt~文件，降级为复制全文到剪贴板\\
~*/\\
\mbox{}\\
import~\{~getChatRows,~getConversations,~getActiveConversationId~\}~from~'./storage.js'\\
import~\{~getConversationSettings~\}~from~'./chat.js'\\
\mbox{}\\
function~\_pad(n)~\{\\
~~~~return~String(n).padStart(2,~'0')\\
\}\\
\mbox{}\\
function~\_nowText()~\{\\
~~~~const~d~=~new~Date()\\
~~~~return~`\$\{d.getFullYear()\}-\$\{\_pad(d.getMonth()~+~1)\}-\$\{\_pad(d.getDate())\}~\$\{\_pad(d.getHours())\}:\\
\$\{\_pad(d.getMinutes())\}`\\
\}\\
\mbox{}\\
/**~组装当前会话聊天记录文本（纯函数，供导出与测试）~*/\\
export~function~buildChatExportText()~\{\\
~~~~const~rows~=~getChatRows()\\
~~~~if~(!rows.length)~return~''\\
~~~~const~conv~=~getConversations().find((c)~=>~c.id~===~getActiveConversationId())\\
~~~~const~s~=~getConversationSettings()\\
~~~~const~lines~=~[]\\
~~~~lines.push(`\$\{s.personaName\}~.~\$\{(conv~\&\&~conv.title)~||~'新对话'\}`)\\
~~~~lines.push(`导出时间：\$\{\_nowText()\}`)\\
~~~~lines.push('--------------------------------')\\
~~~~for~(const~r~of~rows)~\{
}

\newpage

{\ttfamily\footnotesize\raggedright
\mbox{}\\
~~~~.free-tip~\{\\
~~~~~~~~font-size:~22rpx;\\
~~~~~~~~color:~var(--c-text-aid);\\
~~~~~~~~line-height:~1.6;\\
~~~~~~~~margin-bottom:~16rpx;\\
~~~~\}\\
~~~~.free-item~\{\\
~~~~~~~~display:~flex;\\
~~~~~~~~align-items:~center;\\
~~~~~~~~justify-content:~space-between;\\
~~~~~~~~padding:~20rpx~24rpx;\\
~~~~~~~~border-radius:~12rpx;\\
~~~~~~~~background:~var(--c-bg);\\
~~~~~~~~margin-bottom:~16rpx;\\
~~~~~~~~border:~2rpx~solid~transparent;\\
~~~~\}\\
~~~~.free-item:last-child~\{\\
~~~~~~~~margin-bottom:~0;\\
~~~~\}\\
~~~~.free-item.on~\{\\
~~~~~~~~border-color:~var(--c-primary);\\
~~~~~~~~background:~var(--c-primary-light);\\
~~~~\}\\
~~~~.free-main~\{\\
~~~~~~~~flex:~1;\\
~~~~~~~~min-width:~0;\\
~~~~~~~~margin-right:~16rpx;\\
~~~~\}\\
~~~~.free-name~\{\\
~~~~~~~~font-size:~28rpx;\\
~~~~~~~~color:~var(--c-text);\\
~~~~~~~~font-weight:~500;\\
~~~~\}\\
~~~~.free-desc~\{\\
~~~~~~~~font-size:~22rpx;\\
~~~~~~~~color:~var(--c-text-aid);\\
~~~~~~~~margin-top:~6rpx;\\
~~~~~~~~line-height:~1.5;\\
~~~~\}\\
~~~~.free-btn~\{\\
~~~~~~~~flex-shrink:~0;\\
~~~~~~~~font-size:~24rpx;\\
~~~~~~~~color:~var(--c-primary);\\
~~~~~~~~border:~1rpx~solid~var(--c-primary);\\
~~~~~~~~border-radius:~26rpx;\\
~~~~~~~~padding:~8rpx~26rpx;\\
~~~~\}\\
~~~~.free-item.on~.free-btn~\{\\
~~~~~~~~color:~\#fff;
}

\newpage

{\ttfamily\footnotesize\raggedright
~~~~~~~~background:~var(--c-primary);\\
~~~~\}\\
\mbox{}\\
~~~~/*~配置教学弹窗~*/\\
~~~~.mask~\{\\
~~~~~~~~position:~fixed;\\
~~~~~~~~inset:~0;\\
~~~~~~~~background:~rgba(0,~0,~0,~0.45);\\
~~~~~~~~z-index:~999;\\
~~~~~~~~display:~flex;\\
~~~~~~~~align-items:~center;\\
~~~~~~~~justify-content:~center;\\
~~~~\}\\
~~~~.tutorial~\{\\
~~~~~~~~width:~640rpx;\\
~~~~~~~~max-height:~80vh;\\
~~~~~~~~background:~var(--c-card);\\
~~~~~~~~border-radius:~24rpx;\\
~~~~~~~~padding:~32rpx;\\
~~~~~~~~display:~flex;\\
~~~~~~~~flex-direction:~column;\\
~~~~\}\\
~~~~.tutorial-title~\{\\
~~~~~~~~font-size:~32rpx;\\
~~~~~~~~font-weight:~600;\\
~~~~~~~~color:~var(--c-text);\\
~~~~~~~~margin-bottom:~20rpx;\\
~~~~\}\\
~~~~.tutorial-scroll~\{\\
~~~~~~~~flex:~1;\\
~~~~~~~~min-height:~0;\\
~~~~~~~~max-height:~60vh;\\
~~~~\}\\
~~~~.t-step~\{\\
~~~~~~~~margin-bottom:~24rpx;\\
~~~~\}\\
~~~~.t-step-title~\{\\
~~~~~~~~font-size:~26rpx;\\
~~~~~~~~color:~var(--c-primary);\\
~~~~~~~~font-weight:~600;\\
~~~~~~~~margin-bottom:~6rpx;\\
~~~~\}\\
~~~~.t-step-body~\{\\
~~~~~~~~font-size:~24rpx;\\
~~~~~~~~color:~var(--c-text-secondary);\\
~~~~~~~~line-height:~1.7;\\
~~~~\}\\
~~~~.tutorial-close~\{\\
~~~~~~~~margin-top:~20rpx;\\
~~~~~~~~height:~80rpx;
}

\newpage

{\ttfamily\footnotesize\raggedright
~~~~~~~~line-height:~80rpx;\\
~~~~~~~~text-align:~center;\\
~~~~~~~~font-size:~28rpx;\\
~~~~~~~~color:~\#fff;\\
~~~~~~~~background:~var(--c-brand-gradient);\\
~~~~~~~~border-radius:~999rpx;\\
~~~~\}\\
</style>\\
\mbox{}\\
//~================~FILE:~pages/settings/appearance.vue~================\\
<template>\\
~~~~<view~class="ss-page"~:class="themeClass">\\
~~~~~~~~<view~class="ss-tipbar">界面外观：包含深色模式、聊天气泡不透明度与聊天背景。更改后点击底部「\\
保存设置」生效（深色模式与背景即时生效）。</view>\\
\mbox{}\\
~~~~~~~~<view~class="ss-card">\\
~~~~~~~~~~~~<view~class="ss-title">深色模式</view>\\
~~~~~~~~~~~~<view~class="ss-row">\\
~~~~~~~~~~~~~~~~<view~class="ss-row-head">\\
~~~~~~~~~~~~~~~~~~~~<text~class="ss-label"~style="margin-bottom:~0">切换深色~/~浅色</text>\\
~~~~~~~~~~~~~~~~~~~~<view~class="ss-btns">\\
~~~~~~~~~~~~~~~~~~~~~~~~<view~class="ss-btn~first"~:class="\{~on:~\_\_theme~===~'light'~\}"~@tap="pickTh\\
eme('light')">浅色</view>\\
~~~~~~~~~~~~~~~~~~~~~~~~<view~class="ss-btn"~:class="\{~on:~\_\_theme~===~'dark'~\}"~@tap="pickTheme('da\\
rk')">深色</view>\\
~~~~~~~~~~~~~~~~~~~~</view>\\
~~~~~~~~~~~~~~~~</view>\\
~~~~~~~~~~~~~~~~<view~class="ss-hint">\\
~~~~~~~~~~~~~~~~~~~~\{\{~\_\_theme~===~'dark'~?~'当前为深色模式'~:~'当前为浅色模式'~\}\}，即时生效\\
~~~~~~~~~~~~~~~~</view>\\
~~~~~~~~~~~~</view>\\
~~~~~~~~</view>\\
\mbox{}\\
~~~~~~~~<view~class="ss-card">\\
~~~~~~~~~~~~<view~class="ss-title">聊天气泡</view>\\
~~~~~~~~~~~~<view~class="ss-row">\\
~~~~~~~~~~~~~~~~<text~class="ss-label">背景不透明度~\{\{~Math.round((s.bubbleOpacity~||~1)~*~100)~\}\}\%<\\
/text>\\
~~~~~~~~~~~~~~~~<slider~:value="s.bubbleOpacity"~:min="0.2"~:max="1"~:step="0.05"~activeColor="var(-\\
-c-primary)"~@change="onBubbleOpacity"~/>\\
~~~~~~~~~~~~~~~~<view~class="ss-hint">数值越低气泡背景越透明、越能看到下方背景图；文字始终清晰清晰。\\
此设置随对话保存。</view>\\
~~~~~~~~~~~~</view>\\
~~~~~~~~</view>\\
\mbox{}\\
~~~~~~~~<view~class="ss-card">\\
~~~~~~~~~~~~<view~class="ss-title">聊天背景</view>\\
~~~~~~~~~~~~<view~v-if="bg"~class="bg-preview"~:style="bgStyle">\\
~~~~~~~~~~~~~~~~<text~class="bg-preview-tip">当前背景</text>\\
~~~~~~~~~~~~</view>
}

\newpage

{\ttfamily\footnotesize\raggedright
~~~~~~~~~~~~<view~class="bg-actions">\\
~~~~~~~~~~~~~~~~<button~class="bg-btn~primary"~@tap="chooseBg">选择图片</button>\\
~~~~~~~~~~~~~~~~<button~v-if="bg"~class="bg-btn"~@tap="removeBg">移除背景</button>\\
~~~~~~~~~~~~</view>\\
~~~~~~~~~~~~<view~class="ss-hint">选择后即时应用到聊天页并覆盖背景，立即生效、无需保存。</view>\\
~~~~~~~~</view>\\
\mbox{}\\
~~~~~~~~<view~class="ss-save"~@tap="save">保存设置</view>\\
~~~~</view>\\
</template>\\
\mbox{}\\
<script>\\
~~~~import~\{~getConversationSettings,~saveSettings~\}~from~'../../utils/chat.js'\\
~~~~import~\{~getBackgroundImage,~saveBackgroundImage,~removeBackgroundImage~\}~from~'../../utils/stor\\
age.js'\\
~~~~import~\{~setTheme~\}~from~'../../utils/theme.js'\\
~~~~import~\{~addLog~\}~from~'../../utils/log.js'\\
\mbox{}\\
~~~~export~default~\{\\
~~~~~~~~data()~\{\\
~~~~~~~~~~~~return~\{\\
~~~~~~~~~~~~~~~~s:~getConversationSettings(),\\
~~~~~~~~~~~~~~~~bg:~getBackgroundImage()\\
~~~~~~~~~~~~\}\\
~~~~~~~~\},\\
~~~~~~~~computed:~\{\\
~~~~~~~~~~~~bgStyle()~\{\\
~~~~~~~~~~~~~~~~return~this.bg\\
~~~~~~~~~~~~~~~~~~~~?~\{\\
~~~~~~~~~~~~~~~~~~~~~~~~~~~~backgroundImage:~"url('"~+~this.bg~+~"')",\\
~~~~~~~~~~~~~~~~~~~~~~~~~~~~backgroundSize:~'cover',\\
~~~~~~~~~~~~~~~~~~~~~~~~~~~~backgroundPosition:~'center',\\
~~~~~~~~~~~~~~~~~~~~~~~~~~~~backgroundRepeat:~'no-repeat'\\
~~~~~~~~~~~~~~~~~~~~~~\}\\
~~~~~~~~~~~~~~~~~~~~:~\{\}\\
~~~~~~~~~~~~\}\\
~~~~~~~~\},\\
~~~~~~~~onShow()~\{\\
~~~~~~~~~~~~this.s~=~getConversationSettings()\\
~~~~~~~~~~~~this.bg~=~getBackgroundImage()\\
~~~~~~~~\},\\
~~~~~~~~methods:~\{\\
~~~~~~~~~~~~pickTheme(t)~\{\\
~~~~~~~~~~~~~~~~setTheme(t)~//~全局即时生效\\
~~~~~~~~~~~~\},\\
~~~~~~~~~~~~onBubbleOpacity(e)~\{\\
~~~~~~~~~~~~~~~~this.s.bubbleOpacity~=~e.detail.value\\
~~~~~~~~~~~~\},\\
~~~~~~~~~~~~save()~\{\\
~~~~~~~~~~~~~~~~saveSettings(this.s)~//~保存气泡不透明度（随会话快照）
}

\newpage

{\ttfamily\footnotesize\raggedright
~~~~~~~~~~~~~~~~addLog('info',~'保存设置',~'界面外观（气泡不透明度）')\\
~~~~~~~~~~~~~~~~uni.showToast(\{~title:~'已保存',~icon:~'success'~\})\\
~~~~~~~~~~~~\},\\
~~~~~~~~~~~~chooseBg()~\{\\
~~~~~~~~~~~~~~~~uni.chooseImage(\{\\
~~~~~~~~~~~~~~~~~~~~count:~1,\\
~~~~~~~~~~~~~~~~~~~~sizeType:~['compressed'],\\
~~~~~~~~~~~~~~~~~~~~success:~async~(res)~=>~\{\\
~~~~~~~~~~~~~~~~~~~~~~~~const~tempPath~=~res.tempFilePaths[0]\\
~~~~~~~~~~~~~~~~~~~~~~~~if~(!tempPath)~return\\
~~~~~~~~~~~~~~~~~~~~~~~~uni.showLoading(\{~title:~'处理中'~\})\\
~~~~~~~~~~~~~~~~~~~~~~~~try~\{\\
~~~~~~~~~~~~~~~~~~~~~~~~~~~~let~path~=~tempPath\\
~~~~~~~~~~~~~~~~~~~~~~~~~~~~if~(typeof~uni.compressImage~===~'function')~\{\\
~~~~~~~~~~~~~~~~~~~~~~~~~~~~~~~~const~compressed~=~await~new~Promise((r)~=>~\{\\
~~~~~~~~~~~~~~~~~~~~~~~~~~~~~~~~~~~~uni.compressImage(\{\\
~~~~~~~~~~~~~~~~~~~~~~~~~~~~~~~~~~~~~~~~src:~tempPath,\\
~~~~~~~~~~~~~~~~~~~~~~~~~~~~~~~~~~~~~~~~quality:~80,\\
~~~~~~~~~~~~~~~~~~~~~~~~~~~~~~~~~~~~~~~~success:~(x)~=>~r(x.tempFilePath),\\
~~~~~~~~~~~~~~~~~~~~~~~~~~~~~~~~~~~~~~~~fail:~()~=>~r(null)\\
~~~~~~~~~~~~~~~~~~~~~~~~~~~~~~~~~~~~\})\\
~~~~~~~~~~~~~~~~~~~~~~~~~~~~~~~~\})\\
~~~~~~~~~~~~~~~~~~~~~~~~~~~~~~~~if~(compressed)~path~=~compressed\\
~~~~~~~~~~~~~~~~~~~~~~~~~~~~\}\\
~~~~~~~~~~~~~~~~~~~~~~~~~~~~this.bg~=~await~saveBackgroundImage(path)\\
~~~~~~~~~~~~~~~~~~~~~~~~~~~~uni.showToast(\{~title:~'已设置背景',~icon:~'success'~\})\\
~~~~~~~~~~~~~~~~~~~~~~~~\}~finally~\{\\
~~~~~~~~~~~~~~~~~~~~~~~~~~~~uni.hideLoading()\\
~~~~~~~~~~~~~~~~~~~~~~~~\}\\
~~~~~~~~~~~~~~~~~~~~\}\\
~~~~~~~~~~~~~~~~\})\\
~~~~~~~~~~~~\},\\
~~~~~~~~~~~~removeBg()~\{\\
~~~~~~~~~~~~~~~~removeBackgroundImage()\\
~~~~~~~~~~~~~~~~this.bg~=~''\\
~~~~~~~~~~~~~~~~uni.showToast(\{~title:~'已移除背景',~icon:~'success'~\})\\
~~~~~~~~~~~~\}\\
~~~~~~~~\}\\
~~~~\}\\
</script>\\
\mbox{}\\
<style>\\
~~~~.bg-preview~\{\\
~~~~~~~~height:~260rpx;\\
~~~~~~~~border-radius:~12rpx;\\
~~~~~~~~margin-bottom:~20rpx;\\
~~~~~~~~display:~flex;\\
~~~~~~~~align-items:~flex-end;\\
~~~~~~~~justify-content:~center;\\
~~~~~~~~overflow:~hidden;
}

\newpage

{\ttfamily\footnotesize\raggedright
~~~~\}\\
\mbox{}\\
~~~~.bg-preview-tip~\{\\
~~~~~~~~font-size:~22rpx;\\
~~~~~~~~color:~\#fff;\\
~~~~~~~~background:~rgba(0,~0,~0,~0.35);\\
~~~~~~~~padding:~6rpx~20rpx;\\
~~~~~~~~border-radius:~20rpx;\\
~~~~~~~~margin-bottom:~16rpx;\\
~~~~\}\\
\mbox{}\\
~~~~.bg-actions~\{\\
~~~~~~~~display:~flex;\\
~~~~\}\\
\mbox{}\\
~~~~.bg-btn~\{\\
~~~~~~~~flex:~1;\\
~~~~~~~~height:~76rpx;\\
~~~~~~~~line-height:~76rpx;\\
~~~~~~~~font-size:~26rpx;\\
~~~~~~~~color:~var(--c-text-secondary);\\
~~~~~~~~background:~var(--c-bg);\\
~~~~~~~~border-radius:~12rpx;\\
~~~~~~~~margin-right:~16rpx;\\
~~~~\}\\
\mbox{}\\
~~~~.bg-btn:last-child~\{\\
~~~~~~~~margin-right:~0;\\
~~~~\}\\
\mbox{}\\
~~~~.bg-btn.primary~\{\\
~~~~~~~~color:~\#fff;\\
~~~~~~~~background:~var(--c-primary);\\
~~~~\}\\
</style>\\
\mbox{}\\
//~================~FILE:~pages/settings/compress.vue~================\\
<template>\\
~~~~<view~class="ss-page"~:class="themeClass">\\
~~~~~~~~<view~class="ss-tipbar">上下文压缩：长对话自动交给~LLM~压缩为概要，减少后续请求~token~消耗。\\
本页设置更改后自动保存。</view>\\
\mbox{}\\
~~~~~~~~<view~class="ss-auto">[v]~更改自动保存</view>\\
\mbox{}\\
~~~~~~~~<view~class="ss-card">\\
~~~~~~~~~~~~<view~class="ss-title">自动压缩间隔</view>\\
~~~~~~~~~~~~<view~class="ss-row">\\
~~~~~~~~~~~~~~~~<view~class="ss-btns">\\
~~~~~~~~~~~~~~~~~~~~<view\\
~~~~~~~~~~~~~~~~~~~~~~~~v-for="o~in~compressOptions"
}

\newpage

{\ttfamily\footnotesize\raggedright
~~~~~~~~~~~~~~~~~~~~~~~~:key="o.value"\\
~~~~~~~~~~~~~~~~~~~~~~~~class="ss-btn~first"\\
~~~~~~~~~~~~~~~~~~~~~~~~:class="\{~on:~s.compressInterval~===~o.value~\}"\\
~~~~~~~~~~~~~~~~~~~~~~~~@tap="setInterval(o.value)"\\
~~~~~~~~~~~~~~~~~~~~>\{\{~o.label~\}\}</view>\\
~~~~~~~~~~~~~~~~</view>\\
~~~~~~~~~~~~~~~~<view~class="ss-hint">\\
~~~~~~~~~~~~~~~~~~~~累计新增消息达到设定条数后，自动把上文交给~LLM~压缩成概要；历史消息仍完整保留可\\
翻阅，也可在聊天页手动压缩。\\
~~~~~~~~~~~~~~~~</view>\\
~~~~~~~~~~~~</view>\\
~~~~~~~~</view>\\
~~~~</view>\\
</template>\\
\mbox{}\\
<script>\\
~~~~import~\{~getConversationSettings,~saveSettings~\}~from~'../../utils/chat.js'\\
\mbox{}\\
~~~~export~default~\{\\
~~~~~~~~data()~\{\\
~~~~~~~~~~~~return~\{\\
~~~~~~~~~~~~~~~~s:~getConversationSettings(),\\
~~~~~~~~~~~~~~~~compressOptions:~[\\
~~~~~~~~~~~~~~~~~~~~\{~label:~'关闭',~value:~0~\},\\
~~~~~~~~~~~~~~~~~~~~\{~label:~'20条',~value:~20~\},\\
~~~~~~~~~~~~~~~~~~~~\{~label:~'30条',~value:~30~\},\\
~~~~~~~~~~~~~~~~~~~~\{~label:~'40条',~value:~40~\},\\
~~~~~~~~~~~~~~~~~~~~\{~label:~'60条',~value:~60~\},\\
~~~~~~~~~~~~~~~~~~~~\{~label:~'80条',~value:~80~\}\\
~~~~~~~~~~~~~~~~]\\
~~~~~~~~~~~~\}\\
~~~~~~~~\},\\
~~~~~~~~onShow()~\{\\
~~~~~~~~~~~~this.s~=~getConversationSettings()\\
~~~~~~~~\},\\
~~~~~~~~methods:~\{\\
~~~~~~~~~~~~setInterval(v)~\{\\
~~~~~~~~~~~~~~~~if~(this.s.compressInterval~===~v)~return\\
~~~~~~~~~~~~~~~~this.s.compressInterval~=~v\\
~~~~~~~~~~~~~~~~saveSettings(this.s)~//~更改后自动保存\\
~~~~~~~~~~~~~~~~uni.showToast(\{~title:~'已保存（'~+~(v~||~'关闭')~+~'）',~icon:~'none'~\})\\
~~~~~~~~~~~~\}\\
~~~~~~~~\}\\
~~~~\}\\
</script>\\
\mbox{}\\
//~================~FILE:~pages/settings/data.vue~================\\
<template>\\
~~~~<view~class="ss-page"~:class="themeClass">\\
~~~~~~~~<view~class="ss-tipbar">数据管理：清空操作即时生效、无法撤销，请谨慎操作。</view>
}

\newpage

{\ttfamily\footnotesize\raggedright
\mbox{}\\
~~~~~~~~<view~class="ss-card">\\
~~~~~~~~~~~~<view~class="ss-title">数据管理</view>\\
~~~~~~~~~~~~<button~class="ss-danger-btn"~@tap="clearChats">清空对话</button>\\
~~~~~~~~~~~~<button~class="ss-danger-btn"~@tap="clearData">清空记忆与对话</button>\\
~~~~~~~~~~~~<view~class="ss-hint">清空对话仅清空当前会话消息；清空记忆与对话会删除所有记忆与对话记录\\
。小程序端需在公众平台配置~request~合法域名；App~端保持联网即可。</view>\\
~~~~~~~~</view>\\
~~~~</view>\\
</template>\\
\mbox{}\\
<script>\\
~~~~import~\{~clearConversation~\}~from~'../../utils/chat.js'\\
~~~~import~\{~clearAllData~\}~from~'../../utils/storage.js'\\
\mbox{}\\
~~~~export~default~\{\\
~~~~~~~~methods:~\{\\
~~~~~~~~~~~~clearChats()~\{\\
~~~~~~~~~~~~~~~~uni.showModal(\{\\
~~~~~~~~~~~~~~~~~~~~title:~'清空对话',\\
~~~~~~~~~~~~~~~~~~~~content:~'确定清空对话历史吗？',\\
~~~~~~~~~~~~~~~~~~~~success:~(res)~=>~\{\\
~~~~~~~~~~~~~~~~~~~~~~~~if~(res.confirm)~\{\\
~~~~~~~~~~~~~~~~~~~~~~~~~~~~clearConversation()\\
~~~~~~~~~~~~~~~~~~~~~~~~~~~~uni.showToast(\{~title:~'已清空',~icon:~'success'~\})\\
~~~~~~~~~~~~~~~~~~~~~~~~\}\\
~~~~~~~~~~~~~~~~~~~~\}\\
~~~~~~~~~~~~~~~~\})\\
~~~~~~~~~~~~\},\\
~~~~~~~~~~~~clearData()~\{\\
~~~~~~~~~~~~~~~~uni.showModal(\{\\
~~~~~~~~~~~~~~~~~~~~title:~'清空全部数据',\\
~~~~~~~~~~~~~~~~~~~~content:~'将删除所有记忆与对话记录，且不可恢复。确定继续？',\\
~~~~~~~~~~~~~~~~~~~~success:~(res)~=>~\{\\
~~~~~~~~~~~~~~~~~~~~~~~~if~(res.confirm)~\{\\
~~~~~~~~~~~~~~~~~~~~~~~~~~~~clearAllData()\\
~~~~~~~~~~~~~~~~~~~~~~~~~~~~uni.showToast(\{~title:~'已清空',~icon:~'success'~\})\\
~~~~~~~~~~~~~~~~~~~~~~~~\}\\
~~~~~~~~~~~~~~~~~~~~\}\\
~~~~~~~~~~~~~~~~\})\\
~~~~~~~~~~~~\}\\
~~~~~~~~\}\\
~~~~\}\\
</script>\\
\mbox{}\\
<style>\\
~~~~.ss-danger-btn~\{\\
~~~~~~~~margin-bottom:~20rpx;\\
~~~~~~~~height:~80rpx;\\
~~~~~~~~line-height:~80rpx;
}

\newpage

{\ttfamily\footnotesize\raggedright
~~~~~~~~font-size:~28rpx;\\
~~~~~~~~color:~var(--c-danger);\\
~~~~~~~~background:~var(--c-card);\\
~~~~~~~~border:~1rpx~solid~var(--c-danger-light);\\
~~~~~~~~border-radius:~12rpx;\\
~~~~\}\\
</style>\\
\mbox{}\\
//~================~FILE:~pages/settings/help.vue~================\\
<template>\\
~~~~<view~class="ss-page"~:class="themeClass">\\
~~~~~~~~<view~class="ss-tipbar">帮助：了解本~App~的各个功能与使用方法。点击条目展开详情。</view>\\
\mbox{}\\
~~~~~~~~<view~class="ss-card">\\
~~~~~~~~~~~~<view\\
~~~~~~~~~~~~~~~~v-for="(item,~i)~in~helps"\\
~~~~~~~~~~~~~~~~:key="item.title"\\
~~~~~~~~~~~~~~~~class="help-item"\\
~~~~~~~~~~~~>\\
~~~~~~~~~~~~~~~~<view~class="help-head"~@tap="toggle(i)">\\
~~~~~~~~~~~~~~~~~~~~<text~class="help-title">\{\{~item.title~\}\}</text>\\
~~~~~~~~~~~~~~~~~~~~<text~class="help-arrow">\{\{~open~===~i~?~'-'~:~'+'~\}\}</text>\\
~~~~~~~~~~~~~~~~</view>\\
~~~~~~~~~~~~~~~~<view~v-if="open~===~i"~class="help-body">\\
~~~~~~~~~~~~~~~~~~~~<text~class="help-text">\{\{~item.body~\}\}</text>\\
~~~~~~~~~~~~~~~~</view>\\
~~~~~~~~~~~~</view>\\
~~~~~~~~</view>\\
~~~~</view>\\
</template>\\
\mbox{}\\
<script>\\
~~~~export~default~\{\\
~~~~~~~~data()~\{\\
~~~~~~~~~~~~return~\{\\
~~~~~~~~~~~~~~~~open:~0,\\
~~~~~~~~~~~~~~~~helps:~[\\
~~~~~~~~~~~~~~~~~~~~\{\\
~~~~~~~~~~~~~~~~~~~~~~~~title:~'开始一次对话',\\
~~~~~~~~~~~~~~~~~~~~~~~~body:~'在聊天页底部输入框打字并发送即开始。首次使用请先到「设置~->~接口配置\\
」填好接口地址与~API~Key，再回聊天页发消息。'\\
~~~~~~~~~~~~~~~~~~~~\},\\
~~~~~~~~~~~~~~~~~~~~\{\\
~~~~~~~~~~~~~~~~~~~~~~~~title:~'记忆如何运作',\\
~~~~~~~~~~~~~~~~~~~~~~~~body:~'AI~会在回复中自动沉淀你的核心事实、约定与近期状态，按重要性分级（L1~\\
核心~/~L2~约定~/~L3~临时）持续衰减维护。可到「记忆」标签页查看、编辑、删除或手动新建记忆，也可多选批\\
量删除。'\\
~~~~~~~~~~~~~~~~~~~~\},\\
~~~~~~~~~~~~~~~~~~~~\{\\
~~~~~~~~~~~~~~~~~~~~~~~~title:~'使用表情包',
}

\newpage

{\ttfamily\footnotesize\raggedright
~~~~~~~~~~~~~~~~~~~~~~~~body:~'到「表情管理」标签页上传图片并逐张命名，聊天时输入框旁的「表情」按钮\\
可插入~\$表情名\$。开启「人格配置~->~聊天表情包」后，AI~也会在你的表情包清单中挑选使用。'\\
~~~~~~~~~~~~~~~~~~~~\},\\
~~~~~~~~~~~~~~~~~~~~\{\\
~~~~~~~~~~~~~~~~~~~~~~~~title:~'情景（Scene）',\\
~~~~~~~~~~~~~~~~~~~~~~~~body:~'聊天页顶部的情景条会结合当前时间推断「你在做什么」。AI~每次回复都会更\\
新情景；点击情景条可查看、修改或清除，并从下拉历史中快速复用。'\\
~~~~~~~~~~~~~~~~~~~~\},\\
~~~~~~~~~~~~~~~~~~~~\{\\
~~~~~~~~~~~~~~~~~~~~~~~~title:~'拟真聊天',\\
~~~~~~~~~~~~~~~~~~~~~~~~body:~'在「设置~->~拟真聊天」开启后，AI~会在随机时间主动给你发消息（仅现实时\\
间模式、且在当前对话生效），每条~<=1~句。可设置主动时段与频率，或自定义倒计时调试。'\\
~~~~~~~~~~~~~~~~~~~~\},\\
~~~~~~~~~~~~~~~~~~~~\{\\
~~~~~~~~~~~~~~~~~~~~~~~~title:~'语音阅读（TTS）',\\
~~~~~~~~~~~~~~~~~~~~~~~~body:~'在「设置~->~语音阅读」填入~TTS~API~Key（默认~Qwen-TTS）并开启后，AI~\\
新回复会自动朗读，表情不朗读、语音不落盘。可用「测试语音接口」验证配置。'\\
~~~~~~~~~~~~~~~~~~~~\},\\
~~~~~~~~~~~~~~~~~~~~\{\\
~~~~~~~~~~~~~~~~~~~~~~~~title:~'上下文压缩',\\
~~~~~~~~~~~~~~~~~~~~~~~~body:~'长对话会自动（或手动）把较早上文交给~LLM~压缩成概要，节省后续请求的~t\\
oken~消耗；原始消息仍完整保留可翻阅。间隔可在「设置~->~上下文压缩」调整。'\\
~~~~~~~~~~~~~~~~~~~~\},\\
~~~~~~~~~~~~~~~~~~~~\{\\
~~~~~~~~~~~~~~~~~~~~~~~~title:~'重新生成与回撤',\\
~~~~~~~~~~~~~~~~~~~~~~~~body:~'AI~回复下方有「重新生成」可让~AI~重答；重新生成会撤销上一轮回复及其产\\
生的情景/记忆变化，保留你的提问。'\\
~~~~~~~~~~~~~~~~~~~~\},\\
~~~~~~~~~~~~~~~~~~~~\{\\
~~~~~~~~~~~~~~~~~~~~~~~~title:~'多会话管理',\\
~~~~~~~~~~~~~~~~~~~~~~~~body:~'聊天页头部「历史」可查看、切换、删除会话，或把当前对话复制为副本；「\\
新对话」开始新会话并复制当前设置。每个会话拥有独立的记忆、情景与设置。'\\
~~~~~~~~~~~~~~~~~~~~\},\\
~~~~~~~~~~~~~~~~~~~~\{\\
~~~~~~~~~~~~~~~~~~~~~~~~title:~'导出聊天记录',\\
~~~~~~~~~~~~~~~~~~~~~~~~body:~'在聊天页「历史~->~导出对话为~.txt」可将当前会话导出为文本：H5~下载、A\\
pp~保存到文档目录、小程序复制全文。'\\
~~~~~~~~~~~~~~~~~~~~\},\\
~~~~~~~~~~~~~~~~~~~~\{\\
~~~~~~~~~~~~~~~~~~~~~~~~title:~'排查问题',\\
~~~~~~~~~~~~~~~~~~~~~~~~body:~'遇到报错或无回复时，到「设置~->~调试日志」查看请求、响应与错误详情，\\
可据此确认接口地址/模型名是否正确，判断返回的是否为流式数据。'\\
~~~~~~~~~~~~~~~~~~~~\}\\
~~~~~~~~~~~~~~~~]\\
~~~~~~~~~~~~\}\\
~~~~~~~~\},\\
~~~~~~~~methods:~\{\\
~~~~~~~~~~~~toggle(i)~\{\\
~~~~~~~~~~~~~~~~this.open~=~this.open~===~i~?~-1~:~i\\
~~~~~~~~~~~~\}
}

\newpage

{\ttfamily\footnotesize\raggedright
~~~~~~~~\}\\
~~~~\}\\
</script>\\
\mbox{}\\
<style>\\
~~~~.help-item~\{\\
~~~~~~~~border-bottom:~1rpx~solid~var(--c-line);\\
~~~~\}\\
~~~~.help-item:last-child~\{\\
~~~~~~~~border-bottom:~none;\\
~~~~\}\\
~~~~.help-head~\{\\
~~~~~~~~display:~flex;\\
~~~~~~~~align-items:~center;\\
~~~~~~~~justify-content:~space-between;\\
~~~~~~~~padding:~26rpx~6rpx;\\
~~~~\}\\
~~~~.help-title~\{\\
~~~~~~~~font-size:~28rpx;\\
~~~~~~~~color:~var(--c-text);\\
~~~~~~~~font-weight:~500;\\
~~~~\}\\
~~~~.help-arrow~\{\\
~~~~~~~~font-size:~32rpx;\\
~~~~~~~~color:~var(--c-text-aid);\\
~~~~~~~~line-height:~1;\\
~~~~\}\\
~~~~.help-body~\{\\
~~~~~~~~padding:~0~6rpx~24rpx;\\
~~~~\}\\
~~~~.help-text~\{\\
~~~~~~~~font-size:~24rpx;\\
~~~~~~~~color:~var(--c-text-secondary);\\
~~~~~~~~line-height:~1.7;\\
~~~~\}\\
</style>\\
\mbox{}\\
//~================~FILE:~pages/settings/logs.vue~================\\
<template>\\
~~~~<view~class="ss-page"~:class="themeClass">\\
~~~~~~~~<view~class="ss-card">\\
~~~~~~~~~~~~<view~class="ss-title">调试日志</view>\\
~~~~~~~~~~~~<view~class="log-actions">\\
~~~~~~~~~~~~~~~~<button~class="log-btn~primary"~@tap="refreshLogs">刷新</button>\\
~~~~~~~~~~~~~~~~<button~class="log-btn"~@tap="doClearLogs">清空日志</button>\\
~~~~~~~~~~~~~~~~<text~class="log-count">共~\{\{~logs.length~\}\}~条</text>\\
~~~~~~~~~~~~</view>\\
~~~~~~~~~~~~<scroll-view~scroll-y~class="log-list">\\
~~~~~~~~~~~~~~~~<view~v-if="!logs.length"~class="log-empty">暂无日志，进行对话或操作后会自动记录</vi\\
ew>
}

\newpage

{\ttfamily\footnotesize\raggedright
~~~~~~~~~~~~~~~~<view\\
~~~~~~~~~~~~~~~~~~~~v-for="l~in~logs"\\
~~~~~~~~~~~~~~~~~~~~:key="l.id"\\
~~~~~~~~~~~~~~~~~~~~class="log-item"\\
~~~~~~~~~~~~~~~~~~~~:class="\{~expanded:~expandedId~===~l.id~\}"\\
~~~~~~~~~~~~~~~~~~~~@tap="toggleLog(l.id)"\\
~~~~~~~~~~~~~~~~>\\
~~~~~~~~~~~~~~~~~~~~<view~class="log-head">\\
~~~~~~~~~~~~~~~~~~~~~~~~<text~class="log-badge"~:class="'t-'~+~l.type">\{\{~typeName(l.type)~\}\}</text>\\
~~~~~~~~~~~~~~~~~~~~~~~~<text~class="log-msg">\{\{~l.msg~\}\}</text>\\
~~~~~~~~~~~~~~~~~~~~~~~~<text~class="log-time">\{\{~fmtTime(l.time)~\}\}</text>\\
~~~~~~~~~~~~~~~~~~~~</view>\\
~~~~~~~~~~~~~~~~~~~~<view~v-if="l.detail"~class="log-detail">\\
~~~~~~~~~~~~~~~~~~~~~~~~<text~class="log-detail-text">\{\{~expandedId~===~l.id~?~l.detail~:~clipDetail\\
(l.detail)~\}\}</text>\\
~~~~~~~~~~~~~~~~~~~~~~~~<text~class="log-toggle">\{\{~expandedId~===~l.id~?~'收起'~:~'展开详情'~\}\}</te\\
xt>\\
~~~~~~~~~~~~~~~~~~~~</view>\\
~~~~~~~~~~~~~~~~</view>\\
~~~~~~~~~~~~</scroll-view>\\
~~~~~~~~</view>\\
~~~~</view>\\
</template>\\
\mbox{}\\
<script>\\
~~~~import~\{~getLogs,~clearLogs~as~clearDebugLogs~\}~from~'../../utils/log.js'\\
\mbox{}\\
~~~~const~TYPE\_NAMES~=~\{~req:~'请求',~res:~'响应',~err:~'错误',~info:~'信息'~\}\\
\mbox{}\\
~~~~export~default~\{\\
~~~~~~~~data()~\{\\
~~~~~~~~~~~~return~\{\\
~~~~~~~~~~~~~~~~logs:~getLogs(),\\
~~~~~~~~~~~~~~~~expandedId:~''\\
~~~~~~~~~~~~\}\\
~~~~~~~~\},\\
~~~~~~~~onShow()~\{\\
~~~~~~~~~~~~this.logs~=~getLogs()\\
~~~~~~~~\},\\
~~~~~~~~methods:~\{\\
~~~~~~~~~~~~refreshLogs()~\{\\
~~~~~~~~~~~~~~~~this.logs~=~getLogs()\\
~~~~~~~~~~~~~~~~this.expandedId~=~''\\
~~~~~~~~~~~~\},\\
~~~~~~~~~~~~toggleLog(id)~\{\\
~~~~~~~~~~~~~~~~this.expandedId~=~this.expandedId~===~id~?~''~:~id\\
~~~~~~~~~~~~\},\\
~~~~~~~~~~~~clipDetail(d)~\{\\
~~~~~~~~~~~~~~~~const~t~=~String(d~||~'')\\
~~~~~~~~~~~~~~~~return~t.length~>~200~?~t.slice(0,~200)~+~'...'~:~t
}

\newpage

{\ttfamily\footnotesize\raggedright
~~~~~~~~~~~~\},\\
~~~~~~~~~~~~doClearLogs()~\{\\
~~~~~~~~~~~~~~~~uni.showModal(\{\\
~~~~~~~~~~~~~~~~~~~~title:~'清空日志',\\
~~~~~~~~~~~~~~~~~~~~content:~'确定清空全部调试日志吗？',\\
~~~~~~~~~~~~~~~~~~~~success:~(res)~=>~\{\\
~~~~~~~~~~~~~~~~~~~~~~~~if~(!res.confirm)~return\\
~~~~~~~~~~~~~~~~~~~~~~~~clearDebugLogs()\\
~~~~~~~~~~~~~~~~~~~~~~~~this.logs~=~[]\\
~~~~~~~~~~~~~~~~~~~~~~~~this.expandedId~=~''\\
~~~~~~~~~~~~~~~~~~~~~~~~uni.showToast(\{~title:~'已清空日志',~icon:~'success'~\})\\
~~~~~~~~~~~~~~~~~~~~\}\\
~~~~~~~~~~~~~~~~\})\\
~~~~~~~~~~~~\},\\
~~~~~~~~~~~~typeName(t)~\{\\
~~~~~~~~~~~~~~~~return~TYPE\_NAMES[t]~||~t\\
~~~~~~~~~~~~\},\\
~~~~~~~~~~~~fmtTime(iso)~\{\\
~~~~~~~~~~~~~~~~if~(!iso)~return~''\\
~~~~~~~~~~~~~~~~const~d~=~new~Date(iso)\\
~~~~~~~~~~~~~~~~if~(Number.isNaN(d.getTime()))~return~''\\
~~~~~~~~~~~~~~~~const~pad~=~(n)~=>~String(n).padStart(2,~'0')\\
~~~~~~~~~~~~~~~~return~`\$\{pad(d.getMonth()~+~1)\}-\$\{pad(d.getDate())\}~\$\{pad(d.getHours())\}:\$\{pad(d.ge\\
tMinutes())\}:\$\{pad(d.getSeconds())\}`\\
~~~~~~~~~~~~\}\\
~~~~~~~~\}\\
~~~~\}\\
</script>\\
\mbox{}\\
<style>\\
~~~~.log-actions~\{\\
~~~~~~~~display:~flex;\\
~~~~~~~~align-items:~center;\\
~~~~~~~~margin-bottom:~20rpx;\\
~~~~\}\\
~~~~.log-btn~\{\\
~~~~~~~~flex:~none;\\
~~~~~~~~width:~200rpx;\\
~~~~~~~~height:~76rpx;\\
~~~~~~~~line-height:~76rpx;\\
~~~~~~~~font-size:~26rpx;\\
~~~~~~~~color:~var(--c-text-secondary);\\
~~~~~~~~background:~var(--c-bg);\\
~~~~~~~~border-radius:~12rpx;\\
~~~~~~~~margin-right:~16rpx;\\
~~~~\}\\
~~~~.log-btn.primary~\{\\
~~~~~~~~color:~\#fff;\\
~~~~~~~~background:~var(--c-primary);\\
~~~~\}
}

\newpage

{\ttfamily\footnotesize\raggedright
~~~~.log-count~\{\\
~~~~~~~~font-size:~22rpx;\\
~~~~~~~~color:~var(--c-text-aid);\\
~~~~\}\\
~~~~.log-list~\{\\
~~~~~~~~height:~1000rpx;\\
~~~~~~~~background:~var(--c-bg);\\
~~~~~~~~border-radius:~12rpx;\\
~~~~~~~~padding:~8rpx~20rpx;\\
~~~~~~~~box-sizing:~border-box;\\
~~~~\}\\
~~~~.log-empty~\{\\
~~~~~~~~padding:~80rpx~0;\\
~~~~~~~~text-align:~center;\\
~~~~~~~~font-size:~24rpx;\\
~~~~~~~~color:~var(--c-text-aid);\\
~~~~\}\\
~~~~.log-item~\{\\
~~~~~~~~padding:~16rpx~0;\\
~~~~~~~~border-bottom:~1rpx~solid~var(--c-line);\\
~~~~\}\\
~~~~.log-item:last-child~\{\\
~~~~~~~~border-bottom:~none;\\
~~~~\}\\
~~~~.log-item.expanded~\{\\
~~~~~~~~background:~var(--c-card);\\
~~~~~~~~border-radius:~8rpx;\\
~~~~~~~~padding-left:~12rpx;\\
~~~~~~~~padding-right:~12rpx;\\
~~~~\}\\
~~~~.log-head~\{\\
~~~~~~~~display:~flex;\\
~~~~~~~~align-items:~center;\\
~~~~\}\\
~~~~.log-badge~\{\\
~~~~~~~~flex-shrink:~0;\\
~~~~~~~~font-size:~20rpx;\\
~~~~~~~~padding:~2rpx~12rpx;\\
~~~~~~~~border-radius:~16rpx;\\
~~~~~~~~margin-right:~12rpx;\\
~~~~~~~~color:~\#fff;\\
~~~~\}\\
~~~~.log-badge.t-req~\{~background:~var(--c-primary);~\}\\
~~~~.log-badge.t-res~\{~background:~\#00b578;~\}\\
~~~~.log-badge.t-err~\{~background:~var(--c-danger);~\}\\
~~~~.log-badge.t-info~\{~background:~\#86909c;~\}\\
~~~~.log-msg~\{\\
~~~~~~~~flex:~1;\\
~~~~~~~~min-width:~0;\\
~~~~~~~~font-size:~26rpx;
}

\newpage

{\ttfamily\footnotesize\raggedright
~~~~~~~~color:~var(--c-text);\\
~~~~~~~~overflow:~hidden;\\
~~~~~~~~text-overflow:~ellipsis;\\
~~~~~~~~white-space:~nowrap;\\
~~~~\}\\
~~~~.log-time~\{\\
~~~~~~~~flex-shrink:~0;\\
~~~~~~~~margin-left:~12rpx;\\
~~~~~~~~font-size:~20rpx;\\
~~~~~~~~color:~var(--c-text-aid);\\
~~~~\}\\
~~~~.log-detail~\{\\
~~~~~~~~margin-top:~8rpx;\\
~~~~~~~~padding:~12rpx~16rpx;\\
~~~~~~~~background:~var(--c-bg);\\
~~~~~~~~border-radius:~8rpx;\\
~~~~\}\\
~~~~.log-detail-text~\{\\
~~~~~~~~font-size:~22rpx;\\
~~~~~~~~color:~var(--c-text-secondary);\\
~~~~~~~~line-height:~1.5;\\
~~~~~~~~word-break:~break-all;\\
~~~~~~~~white-space:~pre-wrap;\\
~~~~\}\\
~~~~.log-toggle~\{\\
~~~~~~~~display:~inline-block;\\
~~~~~~~~margin-top:~8rpx;\\
~~~~~~~~font-size:~22rpx;\\
~~~~~~~~color:~var(--c-primary);\\
~~~~\}\\
</style>\\
\mbox{}\\
//~================~FILE:~pages/settings/persona.vue~================\\
<template>\\
~~~~<view~class="ss-page"~:class="themeClass">\\
~~~~~~~~<view~class="ss-tipbar">人格配置：选择~AI~扮演的角色，设定情景时间与表情包使用。更改后点击底\\
部「保存设置」生效，仅作用于当前对话。</view>\\
\mbox{}\\
~~~~~~~~<view~class="ss-card">\\
~~~~~~~~~~~~<view~class="ss-title">情景与表情</view>\\
~~~~~~~~~~~~<view~class="ss-row">\\
~~~~~~~~~~~~~~~~<view~class="ss-row-head">\\
~~~~~~~~~~~~~~~~~~~~<text~class="ss-label"~style="margin-bottom:~0">情景时间</text>\\
~~~~~~~~~~~~~~~~~~~~<view~class="ss-btns">\\
~~~~~~~~~~~~~~~~~~~~~~~~<view~class="ss-btn~first"~:class="\{~on:~s.timeMode~===~'real'~\}"~@tap="s.ti\\
meMode~=~'real'">现实时间</view>\\
~~~~~~~~~~~~~~~~~~~~~~~~<view~class="ss-btn"~:class="\{~on:~s.timeMode~===~'virtual'~\}"~@tap="s.timeM\\
ode~=~'virtual'">虚拟时间</view>\\
~~~~~~~~~~~~~~~~~~~~</view>\\
~~~~~~~~~~~~~~~~</view>
}

\newpage

{\ttfamily\footnotesize\raggedright
~~~~~~~~~~~~~~~~<view~class="ss-hint">\\
~~~~~~~~~~~~~~~~~~~~\{\{~s.timeMode~===~'real'~?~'会发送当前真实时间，LLM~据此维护现实情景'~:~'不发送\\
真实时间，情景由~LLM~自由想象（适合角色扮演）'~\}\}\\
~~~~~~~~~~~~~~~~</view>\\
~~~~~~~~~~~~</view>\\
~~~~~~~~~~~~<view~class="ss-row">\\
~~~~~~~~~~~~~~~~<view~class="ss-row-head">\\
~~~~~~~~~~~~~~~~~~~~<text~class="ss-label"~style="margin-bottom:~0">聊天表情包</text>\\
~~~~~~~~~~~~~~~~~~~~<view~class="ss-btns">\\
~~~~~~~~~~~~~~~~~~~~~~~~<view~class="ss-btn~first"~:class="\{~on:~s.emojiEnabled~\}"~@tap="s.emojiEnab\\
led~=~true">允许</view>\\
~~~~~~~~~~~~~~~~~~~~~~~~<view~class="ss-btn"~:class="\{~on:~!s.emojiEnabled~\}"~@tap="s.emojiEnabled~=\\
~false">禁用</view>\\
~~~~~~~~~~~~~~~~~~~~</view>\\
~~~~~~~~~~~~~~~~</view>\\
~~~~~~~~~~~~~~~~<view~class="ss-hint">\\
~~~~~~~~~~~~~~~~~~~~允许时请求携带表情清单，LLM~回复可穿插使用你的表情包；禁用后~LLM~不会主动使用（\\
手动插入与渲染不受影响）。\\
~~~~~~~~~~~~~~~~</view>\\
~~~~~~~~~~~~</view>\\
~~~~~~~~</view>\\
\mbox{}\\
~~~~~~~~<view~class="ss-card">\\
~~~~~~~~~~~~<view~class="ss-title">人格选择</view>\\
~~~~~~~~~~~~<view\\
~~~~~~~~~~~~~~~~v-for="p~in~personalities"\\
~~~~~~~~~~~~~~~~:key="p.id"\\
~~~~~~~~~~~~~~~~class="persona"\\
~~~~~~~~~~~~~~~~:class="\{~active:~s.personalityId~===~p.id~\}"\\
~~~~~~~~~~~~~~~~@tap="selectPersona(p.id)"\\
~~~~~~~~~~~~>\\
~~~~~~~~~~~~~~~~<view~class="persona-head">\\
~~~~~~~~~~~~~~~~~~~~<view~class="persona-main">\\
~~~~~~~~~~~~~~~~~~~~~~~~<view~class="persona-name">\{\{~p.name~\}\}</view>\\
~~~~~~~~~~~~~~~~~~~~~~~~<view~class="persona-desc">\{\{~p.desc~\}\}</view>\\
~~~~~~~~~~~~~~~~~~~~</view>\\
~~~~~~~~~~~~~~~~~~~~<text~class="persona-expand"~@tap.stop="togglePersonaPrompt(p.id)">\\
~~~~~~~~~~~~~~~~~~~~~~~~\{\{~expandedPromptId~===~p.id~?~'收起提示词'~:~'查看提示词'~\}\}\\
~~~~~~~~~~~~~~~~~~~~</text>\\
~~~~~~~~~~~~~~~~</view>\\
~~~~~~~~~~~~~~~~<view~v-if="expandedPromptId~===~p.id"~class="persona-prompt">\{\{~p.prompt~\}\}</view>\\
~~~~~~~~~~~~</view>\\
~~~~~~~~~~~~<view\\
~~~~~~~~~~~~~~~~class="persona"\\
~~~~~~~~~~~~~~~~:class="\{~active:~s.personalityId~===~'custom'~\}"\\
~~~~~~~~~~~~~~~~@tap="selectPersona('custom')"\\
~~~~~~~~~~~~>\\
~~~~~~~~~~~~~~~~<view~class="persona-head">\\
~~~~~~~~~~~~~~~~~~~~<view~class="persona-main">\\
~~~~~~~~~~~~~~~~~~~~~~~~<view~class="persona-name">自定义</view>
}

\newpage

{\ttfamily\footnotesize\raggedright
~~~~~~~~~~~~~~~~~~~~~~~~<view~class="persona-desc">参考下方示例或预设人格的提示词，手写你的专属设定<\\
/view>\\
~~~~~~~~~~~~~~~~~~~~</view>\\
~~~~~~~~~~~~~~~~~~~~<text~class="persona-expand"~@tap.stop="togglePersonaPrompt('\_\_sample\_\_')">\\
~~~~~~~~~~~~~~~~~~~~~~~~\{\{~expandedPromptId~===~'\_\_sample\_\_'~?~'收起示例'~:~'查看完整示例'~\}\}\\
~~~~~~~~~~~~~~~~~~~~</text>\\
~~~~~~~~~~~~~~~~</view>\\
~~~~~~~~~~~~~~~~<view~v-if="expandedPromptId~===~'\_\_sample\_\_'"~class="persona-prompt">\\
~~~~~~~~~~~~~~~~~~~~<text>\{\{~samplePrompt~\}\}</text>\\
~~~~~~~~~~~~~~~~~~~~<view~class="prompt-fill"~@tap.stop="fillSample">一键填入到下方编辑框</view>\\
~~~~~~~~~~~~~~~~</view>\\
~~~~~~~~~~~~</view>\\
~~~~~~~~~~~~<textarea\\
~~~~~~~~~~~~~~~~v-if="s.personalityId~===~'custom'"\\
~~~~~~~~~~~~~~~~class="custom-box"\\
~~~~~~~~~~~~~~~~v-model="s.customPrompt"\\
~~~~~~~~~~~~~~~~maxlength="5000"\\
~~~~~~~~~~~~~~~~placeholder="在这里输入完整的人格设定提示词..."\\
~~~~~~~~~~~~/>\\
~~~~~~~~~~~~<view~v-if="s.personalityId~===~'custom'"~class="char-count">\{\{~s.customPrompt.length~\}\}\\
~/~5000</view>\\
~~~~~~~~</view>\\
\mbox{}\\
~~~~~~~~<view~class="ss-save"~@tap="save">保存设置</view>\\
~~~~</view>\\
</template>\\
\mbox{}\\
<script>\\
~~~~import~\{~getConversationSettings,~saveSettings~\}~from~'../../utils/chat.js'\\
~~~~import~\{~PERSONALITIES,~CUSTOM\_PROMPT\_SAMPLE~\}~from~'../../utils/prompts.js'\\
~~~~import~\{~addLog~\}~from~'../../utils/log.js'\\
\mbox{}\\
~~~~export~default~\{\\
~~~~~~~~data()~\{\\
~~~~~~~~~~~~return~\{\\
~~~~~~~~~~~~~~~~s:~getConversationSettings(),\\
~~~~~~~~~~~~~~~~personalities:~PERSONALITIES,\\
~~~~~~~~~~~~~~~~expandedPromptId:~'',\\
~~~~~~~~~~~~~~~~samplePrompt:~CUSTOM\_PROMPT\_SAMPLE\\
~~~~~~~~~~~~\}\\
~~~~~~~~\},\\
~~~~~~~~onShow()~\{\\
~~~~~~~~~~~~this.s~=~getConversationSettings()\\
~~~~~~~~\},\\
~~~~~~~~methods:~\{\\
~~~~~~~~~~~~selectPersona(id)~\{\\
~~~~~~~~~~~~~~~~this.s.personalityId~=~id\\
~~~~~~~~~~~~\},\\
~~~~~~~~~~~~togglePersonaPrompt(id)~\{\\
~~~~~~~~~~~~~~~~this.expandedPromptId~=~this.expandedPromptId~===~id~?~''~:~id
}

\newpage

{\ttfamily\footnotesize\raggedright
~~~~~~~~~~~~\},\\
~~~~~~~~~~~~fillSample()~\{\\
~~~~~~~~~~~~~~~~this.s.customPrompt~=~CUSTOM\_PROMPT\_SAMPLE\\
~~~~~~~~~~~~~~~~uni.showToast(\{~title:~'已填入示例，可在此基础上修改',~icon:~'none'~\})\\
~~~~~~~~~~~~\},\\
~~~~~~~~~~~~save()~\{\\
~~~~~~~~~~~~~~~~saveSettings(this.s)\\
~~~~~~~~~~~~~~~~addLog('info',~'保存设置',~`人格=\$\{this.s.personalityId\}~.~时间模式=\$\{this.s.timeMod\\
e~===~'virtual'~?~'虚拟'~:~'现实'\}~.~表情包=\$\{this.s.emojiEnabled~?~'允许'~:~'禁用'\}`)\\
~~~~~~~~~~~~~~~~uni.showToast(\{~title:~'已保存',~icon:~'success'~\})\\
~~~~~~~~~~~~\}\\
~~~~~~~~\}\\
~~~~\}\\
</script>\\
\mbox{}\\
<style>\\
~~~~.persona~\{\\
~~~~~~~~padding:~20rpx~24rpx;\\
~~~~~~~~border-radius:~12rpx;\\
~~~~~~~~background:~var(--c-bg);\\
~~~~~~~~margin-bottom:~16rpx;\\
~~~~~~~~border:~2rpx~solid~transparent;\\
~~~~\}\\
~~~~.persona.active~\{\\
~~~~~~~~border-color:~var(--c-primary);\\
~~~~~~~~background:~var(--c-primary-light);\\
~~~~\}\\
~~~~.persona-name~\{\\
~~~~~~~~font-size:~28rpx;\\
~~~~~~~~color:~var(--c-text);\\
~~~~~~~~font-weight:~500;\\
~~~~\}\\
~~~~.persona-desc~\{\\
~~~~~~~~font-size:~22rpx;\\
~~~~~~~~color:~var(--c-text-aid);\\
~~~~~~~~margin-top:~6rpx;\\
~~~~\}\\
~~~~.persona-head~\{\\
~~~~~~~~display:~flex;\\
~~~~~~~~align-items:~flex-start;\\
~~~~\}\\
~~~~.persona-main~\{\\
~~~~~~~~flex:~1;\\
~~~~~~~~min-width:~0;\\
~~~~\}\\
~~~~.persona-expand~\{\\
~~~~~~~~flex-shrink:~0;\\
~~~~~~~~margin-left:~16rpx;\\
~~~~~~~~padding-top:~4rpx;\\
~~~~~~~~font-size:~22rpx;
}

\newpage

{\ttfamily\footnotesize\raggedright
~~~~~~~~color:~var(--c-primary);\\
~~~~\}\\
~~~~.persona-prompt~\{\\
~~~~~~~~margin-top:~16rpx;\\
~~~~~~~~padding:~20rpx;\\
~~~~~~~~background:~var(--c-bg);\\
~~~~~~~~border-radius:~12rpx;\\
~~~~~~~~border:~1rpx~solid~var(--c-line);\\
~~~~~~~~font-size:~24rpx;\\
~~~~~~~~color:~var(--c-text-secondary);\\
~~~~~~~~line-height:~1.7;\\
~~~~~~~~white-space:~pre-wrap;\\
~~~~~~~~word-break:~break-word;\\
~~~~\}\\
~~~~.prompt-fill~\{\\
~~~~~~~~margin-top:~16rpx;\\
~~~~~~~~display:~inline-block;\\
~~~~~~~~font-size:~24rpx;\\
~~~~~~~~color:~var(--c-primary);\\
~~~~~~~~border:~1rpx~solid~var(--c-primary);\\
~~~~~~~~border-radius:~24rpx;\\
~~~~~~~~padding:~8rpx~24rpx;\\
~~~~\}\\
~~~~.custom-box~\{\\
~~~~~~~~width:~100\%;\\
~~~~~~~~height:~800rpx;\\
~~~~~~~~background:~var(--c-bg);\\
~~~~~~~~border-radius:~12rpx;\\
~~~~~~~~padding:~20rpx;\\
~~~~~~~~box-sizing:~border-box;\\
~~~~~~~~font-size:~26rpx;\\
~~~~~~~~color:~var(--c-text);\\
~~~~~~~~line-height:~1.6;\\
~~~~\}\\
~~~~.char-count~\{\\
~~~~~~~~margin-top:~10rpx;\\
~~~~~~~~text-align:~right;\\
~~~~~~~~font-size:~22rpx;\\
~~~~~~~~color:~var(--c-text-aid);\\
~~~~\}\\
</style>\\
\mbox{}\\
//~================~FILE:~pages/settings/proactive.vue~================\\
<template>\\
~~~~<view~class="ss-page"~:class="themeClass">\\
~~~~~~~~<view~class="ss-tipbar">拟真聊天：让~AI~更像真人、在随机时间主动给你发消息。更改后点击底部「\\
保存设置」生效。仅「现实时间」模式生效。若要后台接收通知，请在手机设置打开本应用的通知权限</view>\\
\mbox{}\\
~~~~~~~~<view~class="ss-card">\\
~~~~~~~~~~~~<view~class="ss-title">拟真聊天</view>
}

\newpage

{\ttfamily\footnotesize\raggedright
~~~~~~~~~~~~<view~class="ss-row">\\
~~~~~~~~~~~~~~~~<view~class="ss-row-head">\\
~~~~~~~~~~~~~~~~~~~~<text~class="ss-label"~style="margin-bottom:~0">开关</text>\\
~~~~~~~~~~~~~~~~~~~~<view~class="ss-btns">\\
~~~~~~~~~~~~~~~~~~~~~~~~<view~class="ss-btn~first"~:class="\{~on:~s.proactiveEnabled~\}"~@tap="s.proac\\
tiveEnabled~=~true">开启</view>\\
~~~~~~~~~~~~~~~~~~~~~~~~<view~class="ss-btn"~:class="\{~on:~!s.proactiveEnabled~\}"~@tap="s.proactiveE\\
nabled~=~false">关闭</view>\\
~~~~~~~~~~~~~~~~~~~~</view>\\
~~~~~~~~~~~~~~~~</view>\\
~~~~~~~~~~~~~~~~<view~class="ss-hint">开启后~AI~会在随机时间主动发消息（每条~<=1~句、连续表情~<=2）\\
；开启本对话的所有回复都按该短消息规则约束。</view>\\
~~~~~~~~~~~~</view>\\
~~~~~~~~~~~~<view~class="ss-row">\\
~~~~~~~~~~~~~~~~<view~class="ss-row-head">\\
~~~~~~~~~~~~~~~~~~~~<text~class="ss-label"~style="margin-bottom:~0">主动消息时段</text>\\
~~~~~~~~~~~~~~~~~~~~<view~class="proactive-window">\\
~~~~~~~~~~~~~~~~~~~~~~~~<picker\\
~~~~~~~~~~~~~~~~~~~~~~~~~~~~mode="time"\\
~~~~~~~~~~~~~~~~~~~~~~~~~~~~:value="minToTime(s.proactiveStartMin)"\\
~~~~~~~~~~~~~~~~~~~~~~~~~~~~@change="(e)~=>~s.proactiveStartMin~=~timeToMin(e.detail.value)"\\
~~~~~~~~~~~~~~~~~~~~~~~~>\\
~~~~~~~~~~~~~~~~~~~~~~~~~~~~<view~class="pick-btn">\{\{~minToTime(s.proactiveStartMin)~\}\}</view>\\
~~~~~~~~~~~~~~~~~~~~~~~~</picker>\\
~~~~~~~~~~~~~~~~~~~~~~~~<text~class="window-sep">--</text>\\
~~~~~~~~~~~~~~~~~~~~~~~~<picker\\
~~~~~~~~~~~~~~~~~~~~~~~~~~~~mode="time"\\
~~~~~~~~~~~~~~~~~~~~~~~~~~~~:value="minToTime(s.proactiveEndMin)"\\
~~~~~~~~~~~~~~~~~~~~~~~~~~~~@change="(e)~=>~s.proactiveEndMin~=~timeToMin(e.detail.value)"\\
~~~~~~~~~~~~~~~~~~~~~~~~>\\
~~~~~~~~~~~~~~~~~~~~~~~~~~~~<view~class="pick-btn">\{\{~minToTime(s.proactiveEndMin)~\}\}</view>\\
~~~~~~~~~~~~~~~~~~~~~~~~</picker>\\
~~~~~~~~~~~~~~~~~~~~</view>\\
~~~~~~~~~~~~~~~~</view>\\
~~~~~~~~~~~~~~~~<view~class="ss-hint">仅在此时间段内~AI~会主动发消息（精确到分钟），其余时间静默。</\\
view>\\
~~~~~~~~~~~~</view>\\
~~~~~~~~~~~~<view~class="ss-row">\\
~~~~~~~~~~~~~~~~<view~class="ss-row-head">\\
~~~~~~~~~~~~~~~~~~~~<text~class="ss-label"~style="margin-bottom:~0">主动消息频率</text>\\
~~~~~~~~~~~~~~~~~~~~<view~class="ss-btns">\\
~~~~~~~~~~~~~~~~~~~~~~~~<view~class="ss-btn~first"~:class="\{~on:~s.proactiveCustomSeconds~<=~0~\&\&~s.\\
proactiveLevel~===~'low'~\}"~@tap="setLevel('low')">低频</view>\\
~~~~~~~~~~~~~~~~~~~~~~~~<view~class="ss-btn"~:class="\{~on:~s.proactiveCustomSeconds~<=~0~\&\&~s.proact\\
iveLevel~===~'medium'~\}"~@tap="setLevel('medium')">中频</view>\\
~~~~~~~~~~~~~~~~~~~~~~~~<view~class="ss-btn"~:class="\{~on:~s.proactiveCustomSeconds~<=~0~\&\&~s.proact\\
iveLevel~===~'high'~\}"~@tap="setLevel('high')">高频</view>\\
~~~~~~~~~~~~~~~~~~~~~~~~<view~class="ss-btn"~:class="\{~on:~s.proactiveCustomSeconds~>~0~\}"~@tap="s.p\\
roactiveCustomSeconds~=~s.proactiveCustomSeconds~>~0~?~s.proactiveCustomSeconds~:~30">自定义</view>\\
~~~~~~~~~~~~~~~~~~~~</view>
}

\newpage

{\ttfamily\footnotesize\raggedright
~~~~~~~~~~~~~~~~</view>\\
~~~~~~~~~~~~~~~~<view~class="ss-hint">低频~45\textasciitilde{}120~分钟｜中频~15\textasciitilde{}45~分钟｜高频~5\textasciitilde{}15~分钟。填秒数则固\\
定按该倒计时触发（便于调试）。</view>\\
~~~~~~~~~~~~~~~~<view~class="ss-field">\\
~~~~~~~~~~~~~~~~~~~~<text~class="ss-label">自定义倒计时（秒）</text>\\
~~~~~~~~~~~~~~~~~~~~<input~class="ss-input"~type="number"~v-model="s.proactiveCustomSeconds"~placeho\\
lder="0~=~使用上方档位"~/>\\
~~~~~~~~~~~~~~~~</view>\\
~~~~~~~~~~~~</view>\\
~~~~~~~~~~~~<view~class="ss-row">\\
~~~~~~~~~~~~~~~~<view~class="ss-row-head">\\
~~~~~~~~~~~~~~~~~~~~<text~class="ss-label"~style="margin-bottom:~0">调试</text>\\
~~~~~~~~~~~~~~~~~~~~<view~class="ss-btns">\\
~~~~~~~~~~~~~~~~~~~~~~~~<view~class="ss-btn~test"~:class="\{~disabled:~debugging~\}"~@tap="debugProact\\
ive">立即发送一条主动消息</view>\\
~~~~~~~~~~~~~~~~~~~~</view>\\
~~~~~~~~~~~~~~~~</view>\\
~~~~~~~~~~~~~~~~<view~class="ss-hint">忽略时段/间隔，按当前表单配置立即请求模型返回一条主动消息（发\\
送前自动保存设置）。</view>\\
~~~~~~~~~~~~~~~~<view~class="countdown">\{\{~countdown~\}\}</view>\\
~~~~~~~~~~~~</view>\\
~~~~~~~~</view>\\
\mbox{}\\
~~~~~~~~<view~class="ss-save"~@tap="save">保存设置</view>\\
~~~~</view>\\
</template>\\
\mbox{}\\
<script>\\
~~~~import~\{~getConversationSettings,~saveSettings~\}~from~'../../utils/chat.js'\\
~~~~import~\{~debugProactiveMessage,~rearmProactive,~getProactiveCountdown~\}~from~'../../utils/chat-p\\
roactive.js'\\
~~~~import~\{~requestNotifyPermission~\}~from~'../../utils/notify.js'\\
~~~~import~\{~addLog~\}~from~'../../utils/log.js'\\
\mbox{}\\
~~~~export~default~\{\\
~~~~~~~~data()~\{\\
~~~~~~~~~~~~return~\{\\
~~~~~~~~~~~~~~~~s:~getConversationSettings(),\\
~~~~~~~~~~~~~~~~debugging:~false,\\
~~~~~~~~~~~~~~~~countdown:~'',\\
~~~~~~~~~~~~~~~~\_countdownTimer:~null\\
~~~~~~~~~~~~\}\\
~~~~~~~~\},\\
~~~~~~~~onShow()~\{\\
~~~~~~~~~~~~this.s~=~getConversationSettings()\\
~~~~~~~~~~~~this.updateCountdown()\\
~~~~~~~~~~~~if~(this.\_countdownTimer)~clearInterval(this.\_countdownTimer)\\
~~~~~~~~~~~~this.\_countdownTimer~=~setInterval(()~=>~this.updateCountdown(),~1000)\\
~~~~~~~~\},\\
~~~~~~~~onHide()~\{
}

\newpage

{\ttfamily\footnotesize\raggedright
~~~~~~~~~~~~if~(this.\_countdownTimer)~\{\\
~~~~~~~~~~~~~~~~clearInterval(this.\_countdownTimer)\\
~~~~~~~~~~~~~~~~this.\_countdownTimer~=~null\\
~~~~~~~~~~~~\}\\
~~~~~~~~\},\\
~~~~~~~~methods:~\{\\
~~~~~~~~~~~~setLevel(lv)~\{\\
~~~~~~~~~~~~~~~~this.s.proactiveCustomSeconds~=~0\\
~~~~~~~~~~~~~~~~this.s.proactiveLevel~=~lv\\
~~~~~~~~~~~~\},\\
~~~~~~~~~~~~//~分钟~->~"HH:MM"（钳制~0-1439，非法归~0）\\
~~~~~~~~~~~~minToTime(min)~\{\\
~~~~~~~~~~~~~~~~const~m~=~Math.max(0,~Math.min(1439,~parseInt(min,~10)~||~0))\\
~~~~~~~~~~~~~~~~const~h~=~Math.floor(m~/~60)\\
~~~~~~~~~~~~~~~~const~mm~=~m~\%~60\\
~~~~~~~~~~~~~~~~return~(h~<~10~?~'0'~+~h~:~''~+~h)~+~':'~+~(mm~<~10~?~'0'~+~mm~:~''~+~mm)\\
~~~~~~~~~~~~\},\\
~~~~~~~~~~~~//~"HH:MM"~->~当天第几分钟（0-1439）\\
~~~~~~~~~~~~timeToMin(str)~\{\\
~~~~~~~~~~~~~~~~const~p~=~String(str~||~'').split(':')\\
~~~~~~~~~~~~~~~~const~h~=~parseInt(p[0],~10)~||~0\\
~~~~~~~~~~~~~~~~const~mm~=~parseInt(p[1],~10)~||~0\\
~~~~~~~~~~~~~~~~return~Math.max(0,~Math.min(1439,~h~*~60~+~mm))\\
~~~~~~~~~~~~\},\\
~~~~~~~~~~~~updateCountdown()~\{\\
~~~~~~~~~~~~~~~~const~ms~=~getProactiveCountdown()\\
~~~~~~~~~~~~~~~~if~(ms~===~null)~\{\\
~~~~~~~~~~~~~~~~~~~~const~s~=~this.s~||~getConversationSettings()\\
~~~~~~~~~~~~~~~~~~~~if~(!s.proactiveEnabled)~this.countdown~=~'未开启拟真聊天'\\
~~~~~~~~~~~~~~~~~~~~else~if~(s.timeMode~!==~'real')~this.countdown~=~'仅现实时间模式生效'\\
~~~~~~~~~~~~~~~~~~~~else~this.countdown~=~'等待调度'\\
~~~~~~~~~~~~~~~~\}~else~if~(ms~===~0)~\{\\
~~~~~~~~~~~~~~~~~~~~this.countdown~=~'即将发送主动消息...'\\
~~~~~~~~~~~~~~~~\}~else~\{\\
~~~~~~~~~~~~~~~~~~~~const~total~=~Math.ceil(ms~/~1000)\\
~~~~~~~~~~~~~~~~~~~~const~h~=~Math.floor(total~/~3600)\\
~~~~~~~~~~~~~~~~~~~~const~m~=~Math.floor((total~\%~3600)~/~60)\\
~~~~~~~~~~~~~~~~~~~~const~sec~=~total~\%~60\\
~~~~~~~~~~~~~~~~~~~~this.countdown~=~h~>~0\\
~~~~~~~~~~~~~~~~~~~~~~~~?~`下次主动消息：\$\{h\}~小时~\$\{m\}~分`\\
~~~~~~~~~~~~~~~~~~~~~~~~:~`下次主动消息：\$\{m\}~分~\$\{String(sec).padStart(2,~'0')\}~秒`\\
~~~~~~~~~~~~~~~~\}\\
~~~~~~~~~~~~\},\\
~~~~~~~~~~~~async~debugProactive()~\{\\
~~~~~~~~~~~~~~~~if~(this.debugging)~return\\
~~~~~~~~~~~~~~~~if~(!this.s.apiKey.trim())~\{\\
~~~~~~~~~~~~~~~~~~~~uni.showToast(\{~title:~'请填写~API~Key',~icon:~'none'~\})\\
~~~~~~~~~~~~~~~~~~~~return\\
~~~~~~~~~~~~~~~~\}\\
~~~~~~~~~~~~~~~~saveSettings(this.s)
}

\newpage

{\ttfamily\footnotesize\raggedright
~~~~~~~~~~~~~~~~rearmProactive()\\
~~~~~~~~~~~~~~~~this.debugging~=~true\\
~~~~~~~~~~~~~~~~try~\{\\
~~~~~~~~~~~~~~~~~~~~const~r~=~await~debugProactiveMessage()\\
~~~~~~~~~~~~~~~~~~~~if~(r~\&\&~r.lines~\&\&~r.lines.length)~\{\\
~~~~~~~~~~~~~~~~~~~~~~~~uni.showToast(\{~title:~'已发送~'~+~r.lines.length~+~'~条',~icon:~'success'~\}\\
)\\
~~~~~~~~~~~~~~~~~~~~\}~else~\{\\
~~~~~~~~~~~~~~~~~~~~~~~~uni.showToast(\{~title:~'发送被跳过（请求进行中或未配置~API）',~icon:~'none'~\\
\})\\
~~~~~~~~~~~~~~~~~~~~\}\\
~~~~~~~~~~~~~~~~\}~catch~(e)~\{\\
~~~~~~~~~~~~~~~~~~~~uni.showToast(\{~title:~(e~\&\&~e.message)~||~'发送失败',~icon:~'none'~\})\\
~~~~~~~~~~~~~~~~\}~finally~\{\\
~~~~~~~~~~~~~~~~~~~~this.debugging~=~false\\
~~~~~~~~~~~~~~~~\}\\
~~~~~~~~~~~~\},\\
~~~~~~~~~~~~save()~\{\\
~~~~~~~~~~~~~~~~if~(!this.s.apiKey.trim())~\{\\
~~~~~~~~~~~~~~~~~~~~uni.showToast(\{~title:~'请填写~API~Key',~icon:~'none'~\})\\
~~~~~~~~~~~~~~~~~~~~return\\
~~~~~~~~~~~~~~~~\}\\
~~~~~~~~~~~~~~~~saveSettings(this.s)\\
~~~~~~~~~~~~~~~~rearmProactive()\\
~~~~~~~~~~~~~~~~//~开启拟真时申请系统通知权限（后台收到主动消息时弹通知）\\
~~~~~~~~~~~~~~~~if~(this.s.proactiveEnabled)~requestNotifyPermission()\\
~~~~~~~~~~~~~~~~addLog(\\
~~~~~~~~~~~~~~~~~~~~'info',\\
~~~~~~~~~~~~~~~~~~~~'保存设置',\\
~~~~~~~~~~~~~~~~~~~~`拟真=\$\{this.s.proactiveEnabled~?~'开(时段'~+~this.minToTime(this.s.proactiveSta\\
rtMin)~+~'-'~+~this.minToTime(this.s.proactiveEndMin)~+~','~+~(this.s.proactiveCustomSeconds~>~0~?~'\\
自定义'~+~this.s.proactiveCustomSeconds~+~'秒'~:~this.s.proactiveLevel~+~'档')~+~')'~:~'关'\}`\\
~~~~~~~~~~~~~~~~)\\
~~~~~~~~~~~~~~~~uni.showToast(\{~title:~'已保存',~icon:~'success'~\})\\
~~~~~~~~~~~~\}\\
~~~~~~~~\}\\
~~~~\}\\
</script>\\
\mbox{}\\
<style>\\
~~~~.proactive-window~\{\\
~~~~~~~~display:~flex;\\
~~~~~~~~align-items:~center;\\
~~~~\}\\
~~~~.pick-btn~\{\\
~~~~~~~~padding:~8rpx~24rpx;\\
~~~~~~~~border-radius:~12rpx;\\
~~~~~~~~background:~var(--c-card);\\
~~~~~~~~border:~1rpx~solid~var(--c-line);\\
~~~~~~~~font-size:~26rpx;
}

\newpage

{\ttfamily\footnotesize\raggedright
~~~~~~~~color:~var(--c-text);\\
~~~~\}\\
~~~~.window-sep~\{\\
~~~~~~~~margin:~0~16rpx;\\
~~~~~~~~color:~var(--c-text-aid);\\
~~~~~~~~font-size:~26rpx;\\
~~~~\}\\
~~~~.countdown~\{\\
~~~~~~~~margin-top:~12rpx;\\
~~~~~~~~color:~var(--c-primary);\\
~~~~~~~~font-weight:~600;\\
~~~~~~~~font-size:~24rpx;\\
~~~~\}\\
</style>\\
\mbox{}\\
//~================~FILE:~pages/settings/settings.vue~================\\
<template>\\
~~~~<view~class="page"~:class="themeClass">\\
~~~~~~~~<view~class="hero">\\
~~~~~~~~~~~~<view~class="hero-title">设置</view>\\
~~~~~~~~~~~~<view~class="hero-sub">每项设置进入对应页面调整，更改后记得保存</view>\\
~~~~~~~~</view>\\
\mbox{}\\
~~~~~~~~<view~class="group"~v-for="g~in~groups"~:key="g.name">\\
~~~~~~~~~~~~<view~class="group-name">\{\{~g.name~\}\}</view>\\
~~~~~~~~~~~~<view~class="group-card">\\
~~~~~~~~~~~~~~~~<view\\
~~~~~~~~~~~~~~~~~~~~v-for="item~in~g.items"\\
~~~~~~~~~~~~~~~~~~~~:key="item.page"\\
~~~~~~~~~~~~~~~~~~~~class="entry"\\
~~~~~~~~~~~~~~~~~~~~@tap="go(item.page)"\\
~~~~~~~~~~~~~~~~>\\
~~~~~~~~~~~~~~~~~~~~<view~class="entry-main">\\
~~~~~~~~~~~~~~~~~~~~~~~~<view~class="entry-name">\\
~~~~~~~~~~~~~~~~~~~~~~~~~~~~\{\{~item.name~\}\}\\
~~~~~~~~~~~~~~~~~~~~~~~~~~~~<text~v-if="item.required~\&\&~needConfig"~class="req">*</text>\\
~~~~~~~~~~~~~~~~~~~~~~~~~~~~<text~v-if="item.tag"~class="tag">\{\{~item.tag~\}\}</text>\\
~~~~~~~~~~~~~~~~~~~~~~~~</view>\\
~~~~~~~~~~~~~~~~~~~~~~~~<view~class="entry-desc">\{\{~item.desc~\}\}</view>\\
~~~~~~~~~~~~~~~~~~~~</view>\\
~~~~~~~~~~~~~~~~~~~~<text~class="entry-arrow">></text>\\
~~~~~~~~~~~~~~~~</view>\\
~~~~~~~~~~~~</view>\\
~~~~~~~~</view>\\
\mbox{}\\
~~~~<!--~自绘底部导航栏（替代原生~tabBar，随主题深浅色切换）~-->\\
~~~~<app-tab-bar~:active="2"~/>\\
</view>\\
</template>\\
\mbox{}
}

\newpage

{\ttfamily\footnotesize\raggedright
<script>\\
~~~~import~\{~getConversationSettings~\}~from~'../../utils/chat.js'\\
~~~~import~AppTabBar~from~'../../components/app-tab-bar/app-tab-bar.vue'\\
\mbox{}\\
~~~~export~default~\{\\
~~~~~~~~components:~\{~AppTabBar~\},\\
~~~~~~~~data()~\{\\
~~~~~~~~~~~~return~\{\\
~~~~~~~~~~~~~~~~needConfig:~false,~//~接口配置缺失（baseUrl/apiKey~未填）时置~true\\
~~~~~~~~~~~~~~~~groups:~[\\
~~~~~~~~~~~~~~~~~~~~\{\\
~~~~~~~~~~~~~~~~~~~~~~~~name:~'常规',\\
~~~~~~~~~~~~~~~~~~~~~~~~items:~[\\
~~~~~~~~~~~~~~~~~~~~~~~~~~~~\{~page:~'/pages/settings/appearance',~name:~'界面外观',~desc:~'深色模式\\
、聊天气泡不透明度与聊天背景',~tag:~'保存'~\},\\
~~~~~~~~~~~~~~~~~~~~~~~~~~~~\{~page:~'/pages/settings/api',~name:~'接口配置',~desc:~'LLM~接口地址~/~K\\
ey~/~模型~/~预设与配置教学',~required:~true,~tag:~'保存'~\},\\
~~~~~~~~~~~~~~~~~~~~~~~~~~~~\{~page:~'/pages/settings/persona',~name:~'人格配置',~desc:~'人格选择、情\\
景时间、自定义设定',~tag:~'保存'~\},\\
~~~~~~~~~~~~~~~~~~~~~~~~~~~~\{~page:~'/pages/settings/proactive',~name:~'拟真聊天',~desc:~'模拟真人主\\
动发消息的时段与频率',~tag:~'保存'~\},\\
~~~~~~~~~~~~~~~~~~~~~~~~~~~~\{~page:~'/pages/settings/tts',~name:~'语音阅读',~desc:~'回复自动朗读的~T\\
TS~接口与音色',~tag:~'保存'~\}\\
~~~~~~~~~~~~~~~~~~~~~~~~]\\
~~~~~~~~~~~~~~~~~~~~\},\\
~~~~~~~~~~~~~~~~~~~~\{\\
~~~~~~~~~~~~~~~~~~~~~~~~name:~'对话与性能',\\
~~~~~~~~~~~~~~~~~~~~~~~~items:~[\\
~~~~~~~~~~~~~~~~~~~~~~~~~~~~\{~page:~'/pages/settings/compress',~name:~'上下文压缩',~desc:~'长对话自\\
动压缩为概要，节省~token',~tag:~'自动保存'~\}\\
~~~~~~~~~~~~~~~~~~~~~~~~]\\
~~~~~~~~~~~~~~~~~~~~\},\\
~~~~~~~~~~~~~~~~~~~~\{\\
~~~~~~~~~~~~~~~~~~~~~~~~name:~'数据与调试',\\
~~~~~~~~~~~~~~~~~~~~~~~~items:~[\\
~~~~~~~~~~~~~~~~~~~~~~~~~~~~\{~page:~'/pages/settings/data',~name:~'数据管理',~desc:~'清空对话~/~清空\\
记忆与全部数据'~\},\\
~~~~~~~~~~~~~~~~~~~~~~~~~~~~\{~page:~'/pages/settings/logs',~name:~'调试日志',~desc:~'查看请求、响应\\
、错误与操作记录'~\}\\
~~~~~~~~~~~~~~~~~~~~~~~~]\\
~~~~~~~~~~~~~~~~~~~~\},\\
~~~~~~~~~~~~~~~~~~~~\{\\
~~~~~~~~~~~~~~~~~~~~~~~~name:~'关于',\\
~~~~~~~~~~~~~~~~~~~~~~~~items:~[\\
~~~~~~~~~~~~~~~~~~~~~~~~~~~~\{~page:~'/pages/settings/help',~name:~'帮助',~desc:~'各功能介绍与使用方\\
法'~\},\\
~~~~~~~~~~~~~~~~~~~~~~~~~~~~\{~page:~'/pages/settings/about',~name:~'关于',~desc:~'App~版本与技术信息\\
'~\}\\
~~~~~~~~~~~~~~~~~~~~~~~~]\\
~~~~~~~~~~~~~~~~~~~~\}
}

\newpage

{\ttfamily\footnotesize\raggedright
~~~~~~~~~~~~~~~~]\\
~~~~~~~~~~~~\}\\
~~~~~~~~\},\\
~~~~~~~~onShow()~\{\\
~~~~~~~~~~~~//~接口配置为必要项：任一缺失（地址或~Key~为空）即显示红色~*~提示\\
~~~~~~~~~~~~const~s~=~getConversationSettings()\\
~~~~~~~~~~~~this.needConfig~=~!(s.baseUrl~\&\&~s.baseUrl.trim()~\&\&~s.apiKey~\&\&~s.apiKey.trim())\\
~~~~~~~~\},\\
~~~~~~~~methods:~\{\\
~~~~~~~~~~~~go(page)~\{\\
~~~~~~~~~~~~~~~~uni.navigateTo(\{~url:~page~\})\\
~~~~~~~~~~~~\}\\
~~~~~~~~\}\\
~~~~\}\\
</script>\\
\mbox{}\\
<style>\\
~~~~.page~\{\\
~~~~~~~~padding:~20rpx~30rpx~calc(var(--ctab-h,~100rpx)~+~env(safe-area-inset-bottom)~+~40rpx);\\
~~~~~~~~background:~var(--c-bg);\\
~~~~~~~~min-height:~100vh;\\
~~~~~~~~box-sizing:~border-box;\\
~~~~\}\\
\mbox{}\\
~~~~.hero~\{\\
~~~~~~~~padding:~20rpx~6rpx~28rpx;\\
~~~~\}\\
~~~~.hero-title~\{\\
~~~~~~~~font-size:~44rpx;\\
~~~~~~~~font-weight:~700;\\
~~~~~~~~color:~var(--c-text);\\
~~~~\}\\
~~~~.hero-sub~\{\\
~~~~~~~~margin-top:~8rpx;\\
~~~~~~~~font-size:~24rpx;\\
~~~~~~~~color:~var(--c-text-aid);\\
~~~~\}\\
\mbox{}\\
~~~~.group~\{\\
~~~~~~~~margin-bottom:~24rpx;\\
~~~~\}\\
~~~~.group-name~\{\\
~~~~~~~~font-size:~24rpx;\\
~~~~~~~~color:~var(--c-text-aid);\\
~~~~~~~~margin:~0~6rpx~12rpx;\\
~~~~\}\\
~~~~.group-card~\{\\
~~~~~~~~background:~var(--c-card);\\
~~~~~~~~border-radius:~var(--c-radius-lg);\\
~~~~~~~~box-shadow:~var(--c-shadow-card);
}

\newpage

{\ttfamily\footnotesize\raggedright
~~~~~~~~overflow:~hidden;\\
~~~~~~~~border:~1rpx~solid~transparent;\\
~~~~\}\\
\mbox{}\\
~~~~.entry~\{\\
~~~~~~~~display:~flex;\\
~~~~~~~~align-items:~center;\\
~~~~~~~~padding:~26rpx~28rpx;\\
~~~~~~~~border-bottom:~1rpx~solid~var(--c-line);\\
~~~~\}\\
~~~~.entry:last-child~\{\\
~~~~~~~~border-bottom:~none;\\
~~~~\}\\
~~~~.entry-name~\{\\
~~~~~~~~font-size:~28rpx;\\
~~~~~~~~color:~var(--c-text);\\
~~~~~~~~font-weight:~500;\\
~~~~~~~~display:~flex;\\
~~~~~~~~align-items:~center;\\
~~~~\}\\
~~~~.entry-name~.req~\{\\
~~~~~~~~color:~var(--c-danger);\\
~~~~~~~~margin-left:~6rpx;\\
~~~~~~~~font-weight:~700;\\
~~~~\}\\
~~~~.entry-name~.tag~\{\\
~~~~~~~~margin-left:~12rpx;\\
~~~~~~~~font-size:~20rpx;\\
~~~~~~~~color:~var(--c-text-aid);\\
~~~~~~~~font-weight:~400;\\
~~~~~~~~border:~1rpx~solid~var(--c-line);\\
~~~~~~~~border-radius:~20rpx;\\
~~~~~~~~padding:~2rpx~12rpx;\\
~~~~\}\\
~~~~.entry-desc~\{\\
~~~~~~~~margin-top:~6rpx;\\
~~~~~~~~font-size:~22rpx;\\
~~~~~~~~color:~var(--c-text-aid);\\
~~~~~~~~line-height:~1.5;\\
~~~~\}\\
~~~~.entry-arrow~\{\\
~~~~~~~~margin-left:~16rpx;\\
~~~~~~~~font-size:~40rpx;\\
~~~~~~~~color:~var(--c-text-aid);\\
~~~~~~~~line-height:~1;\\
~~~~\}\\
</style>\\
\mbox{}\\
//~================~FILE:~pages/settings/tts.vue~================\\
<template>
}

\newpage

{\ttfamily\footnotesize\raggedright
~~~~<view~class="ss-page"~:class="themeClass">\\
~~~~~~~~<view~class="ss-tipbar">语音阅读：LLM~新返回的消息在渲染时自动合成语音朗读，表情不朗读、语音\\
不落盘。更改后点击底部「保存设置」生效。</view>\\
\mbox{}\\
~~~~~~~~<view~class="ss-card">\\
~~~~~~~~~~~~<view~class="ss-title">语音阅读</view>\\
~~~~~~~~~~~~<view~class="ss-row">\\
~~~~~~~~~~~~~~~~<view~class="ss-row-head">\\
~~~~~~~~~~~~~~~~~~~~<text~class="ss-label"~style="margin-bottom:~0">开关</text>\\
~~~~~~~~~~~~~~~~~~~~<view~class="ss-btns">\\
~~~~~~~~~~~~~~~~~~~~~~~~<view~class="ss-btn~first"~:class="\{~on:~s.ttsEnabled~\}"~@tap="s.ttsEnabled~\\
=~true">开启</view>\\
~~~~~~~~~~~~~~~~~~~~~~~~<view~class="ss-btn"~:class="\{~on:~!s.ttsEnabled~\}"~@tap="s.ttsEnabled~=~fal\\
se">关闭</view>\\
~~~~~~~~~~~~~~~~~~~~</view>\\
~~~~~~~~~~~~~~~~</view>\\
~~~~~~~~~~~~~~~~<view~class="ss-hint">默认以~Qwen-TTS~为例。开启后需填下面的~API~Key~才会朗读。</vie\\
w>\\
~~~~~~~~~~~~</view>\\
~~~~~~~~~~~~<view~class="ss-field">\\
~~~~~~~~~~~~~~~~<text~class="ss-label">TTS~API~Key</text>\\
~~~~~~~~~~~~~~~~<input~class="ss-input"~v-model="s.ttsApiKey"~password~placeholder="sk-..."~/>\\
~~~~~~~~~~~~</view>\\
~~~~~~~~~~~~<view~class="ss-field">\\
~~~~~~~~~~~~~~~~<text~class="ss-label">TTS~模型</text>\\
~~~~~~~~~~~~~~~~<input~class="ss-input"~v-model="s.ttsModel"~placeholder="qwen3-tts-flash"~/>\\
~~~~~~~~~~~~</view>\\
~~~~~~~~~~~~<view~class="ss-field">\\
~~~~~~~~~~~~~~~~<text~class="ss-label">音色</text>\\
~~~~~~~~~~~~~~~~<view~class="voice-row">\\
~~~~~~~~~~~~~~~~~~~~<input~class="ss-input~voice-input"~v-model="s.ttsVoice"~placeholder="Cherry"~/>\\
~~~~~~~~~~~~~~~~~~~~<picker~mode="selector"~:range="ttsVoiceNames"~@change="onTtsVoiceChange">\\
~~~~~~~~~~~~~~~~~~~~~~~~<view~class="voice-pick-btn">从列表选择</view>\\
~~~~~~~~~~~~~~~~~~~~</picker>\\
~~~~~~~~~~~~~~~~</view>\\
~~~~~~~~~~~~~~~~<view~v-if="ttsVoiceHint"~class="ss-hint">\{\{~ttsVoiceHint~\}\}</view>\\
~~~~~~~~~~~~~~~~<view~class="ss-hint">可手动输入官方列表外的自定义音色；完整音色列表参考官方文档。</\\
view>\\
~~~~~~~~~~~~</view>\\
~~~~~~~~~~~~<view~class="ss-field">\\
~~~~~~~~~~~~~~~~<view~class="tts-test-btn"~@tap="doTestTts">测试语音接口</view>\\
~~~~~~~~~~~~~~~~<view~class="ss-hint">按上方当前配置（Key/模型/音色）发送一小段测试文本并尝试播放。<\\
/view>\\
~~~~~~~~~~~~</view>\\
~~~~~~~~</view>\\
\mbox{}\\
~~~~~~~~<view~class="ss-save"~@tap="save">保存设置</view>\\
~~~~</view>\\
</template>\\
\mbox{}
}

\newpage

{\ttfamily\footnotesize\raggedright
<script>\\
~~~~import~\{~getConversationSettings,~saveSettings~\}~from~'../../utils/chat.js'\\
~~~~import~\{~TTS\_VOICE\_NAMES,~ttsVoiceDesc,~testTts~as~runTtsTest~\}~from~'../../utils/tts.js'\\
~~~~import~\{~addLog~\}~from~'../../utils/log.js'\\
\mbox{}\\
~~~~export~default~\{\\
~~~~~~~~data()~\{\\
~~~~~~~~~~~~return~\{\\
~~~~~~~~~~~~~~~~s:~getConversationSettings(),\\
~~~~~~~~~~~~~~~~ttsVoiceNames:~TTS\_VOICE\_NAMES\\
~~~~~~~~~~~~\}\\
~~~~~~~~\},\\
~~~~~~~~computed:~\{\\
~~~~~~~~~~~~ttsVoiceHint()~\{\\
~~~~~~~~~~~~~~~~const~desc~=~ttsVoiceDesc(this.s.ttsVoice)\\
~~~~~~~~~~~~~~~~return~desc~?~'「'~+~this.s.ttsVoice~+~'」'~+~desc~:~''\\
~~~~~~~~~~~~\}\\
~~~~~~~~\},\\
~~~~~~~~onShow()~\{\\
~~~~~~~~~~~~this.s~=~getConversationSettings()\\
~~~~~~~~\},\\
~~~~~~~~methods:~\{\\
~~~~~~~~~~~~onTtsVoiceChange(e)~\{\\
~~~~~~~~~~~~~~~~this.s.ttsVoice~=~this.ttsVoiceNames[e.detail.value]\\
~~~~~~~~~~~~\},\\
~~~~~~~~~~~~async~doTestTts()~\{\\
~~~~~~~~~~~~~~~~if~(!this.s.ttsApiKey.trim())~\{\\
~~~~~~~~~~~~~~~~~~~~uni.showToast(\{~title:~'请先填写~TTS~API~Key',~icon:~'none'~\})\\
~~~~~~~~~~~~~~~~~~~~return\\
~~~~~~~~~~~~~~~~\}\\
~~~~~~~~~~~~~~~~uni.showLoading(\{~title:~'测试中'~\})\\
~~~~~~~~~~~~~~~~let~res~=~\{~ok:~false,~message:~'测试失败'~\}\\
~~~~~~~~~~~~~~~~try~\{\\
~~~~~~~~~~~~~~~~~~~~res~=~await~runTtsTest(\{~apiKey:~this.s.ttsApiKey,~model:~this.s.ttsModel,~voice\\
:~this.s.ttsVoice~\})\\
~~~~~~~~~~~~~~~~\}~catch~(e)~\{\\
~~~~~~~~~~~~~~~~~~~~res~=~\{~ok:~false,~message:~(e~\&\&~e.message)~||~'测试失败'~\}\\
~~~~~~~~~~~~~~~~\}~finally~\{\\
~~~~~~~~~~~~~~~~~~~~uni.hideLoading()\\
~~~~~~~~~~~~~~~~\}\\
~~~~~~~~~~~~~~~~uni.showToast(\{~title:~res.message,~icon:~res.ok~?~'success'~:~'none'~\})\\
~~~~~~~~~~~~\},\\
~~~~~~~~~~~~save()~\{\\
~~~~~~~~~~~~~~~~saveSettings(this.s)\\
~~~~~~~~~~~~~~~~addLog('info',~'保存设置',~`语音阅读=\$\{this.s.ttsEnabled~?~'开('~+~(this.s.ttsModel~\\
||~'')~+~'/'~+~(this.s.ttsVoice~||~'')~+~')'~:~'关'\}`)\\
~~~~~~~~~~~~~~~~uni.showToast(\{~title:~'已保存',~icon:~'success'~\})\\
~~~~~~~~~~~~\}\\
~~~~~~~~\}\\
~~~~\}
}

\newpage

{\ttfamily\footnotesize\raggedright
</script>\\
\mbox{}\\
<style>\\
~~~~.voice-row~\{\\
~~~~~~~~display:~flex;\\
~~~~~~~~align-items:~center;\\
~~~~\}\\
~~~~.voice-input~\{\\
~~~~~~~~flex:~1;\\
~~~~~~~~min-width:~0;\\
~~~~~~~~margin-right:~16rpx;\\
~~~~\}\\
~~~~.voice-pick-btn~\{\\
~~~~~~~~flex-shrink:~0;\\
~~~~~~~~height:~80rpx;\\
~~~~~~~~line-height:~80rpx;\\
~~~~~~~~padding:~0~28rpx;\\
~~~~~~~~font-size:~26rpx;\\
~~~~~~~~color:~var(--c-primary);\\
~~~~~~~~background:~var(--c-primary-light);\\
~~~~~~~~border-radius:~12rpx;\\
~~~~\}\\
~~~~.tts-test-btn~\{\\
~~~~~~~~height:~80rpx;\\
~~~~~~~~line-height:~80rpx;\\
~~~~~~~~text-align:~center;\\
~~~~~~~~font-size:~28rpx;\\
~~~~~~~~color:~\#fff;\\
~~~~~~~~background:~var(--c-brand-gradient);\\
~~~~~~~~border-radius:~var(--c-radius-md);\\
~~~~~~~~box-shadow:~0~4rpx~12rpx~rgba(91,~124,~250,~0.3);\\
~~~~\}\\
</style>\\
\mbox{}\\
//~================~FILE:~uni.promisify.adaptor.js~================\\
uni.addInterceptor(\{\\
~~returnValue~(res)~\{\\
~~~~if~(!(!!res~\&\&~(typeof~res~===~"object"~||~typeof~res~===~"function")~\&\&~typeof~res.then~===~"fu\\
nction"))~\{\\
~~~~~~return~res;\\
~~~~\}\\
~~~~return~new~Promise((resolve,~reject)~=>~\{\\
~~~~~~res.then((res)~=>~\{\\
~~~~~~~~if~(!res)~return~resolve(res)~\\
~~~~~~~~return~res[0]~?~reject(res[0])~:~resolve(res[1])\\
~~~~~~\});\\
~~~~\});\\
~~\},\\
\});\\
\mbox{}
}

\end{document}
